\documentclass[11pt]{article}

\usepackage[margin=1in]{geometry}
\usepackage[T1]{fontenc}
\usepackage{lmodern}
\usepackage{microtype}
\usepackage{amsmath,amssymb,amsthm,mathtools}
\usepackage{booktabs,tabularx}
\usepackage{enumitem}
\usepackage{aliascnt}
\usepackage[dvipsnames]{xcolor}
\usepackage[colorlinks=true,linkcolor=MidnightBlue,citecolor=BrickRed,urlcolor=MidnightBlue]{hyperref}
\usepackage[nameinlink,capitalise,noabbrev]{cleveref}

\setlength{\parindent}{0pt}
\setlength{\parskip}{0.6em}
\setlist{itemsep=0.25em,topsep=0.35em}

\newtheorem{theorem}{Theorem}[section]
\newaliascnt{lemma}{theorem}
\newtheorem{lemma}[lemma]{Lemma}
\aliascntresetthe{lemma}
\newaliascnt{proposition}{theorem}
\newtheorem{proposition}[proposition]{Proposition}
\aliascntresetthe{proposition}
\newaliascnt{corollary}{theorem}
\newtheorem{corollary}[corollary]{Corollary}
\aliascntresetthe{corollary}

% arXiv's TeX Live 2025 currently requires explicit cleveref aliases for
% theorem-like environments that share a section counter.
\AddToHook{env/theorem/begin}{\crefalias{section}{theorem}}
\AddToHook{env/lemma/begin}{\crefalias{section}{lemma}}
\AddToHook{env/proposition/begin}{\crefalias{section}{proposition}}
\AddToHook{env/corollary/begin}{\crefalias{section}{corollary}}

\newcommand{\bits}{\{0,1\}}
\newcommand{\F}{\mathbb{F}}
\newcommand{\C}{\mathrm{C}}
\newcommand{\poly}{\operatorname{poly}}
\newcommand{\xor}{\mathbin{\oplus}}
\newcommand{\bigxor}{\bigoplus}

% The checked-in public and arXiv manuscript is identified.  Switch to
% \anonymoustrue only when producing a separate double-blind review copy.
\newif\ifanonymous
\anonymousfalse

\title{\textbf{Exponential-Range Mass Production of Boolean Functions}\\[0.35em]
\large Local recovery and optimal quantum state preparation}
\ifanonymous
  \author{}
  \date{}
  \hypersetup{
    pdftitle={Exponential-Range Mass Production of Boolean Functions: Local recovery and optimal quantum state preparation},
    pdfauthor={}
  }
\else
  \author{Samuel Schlesinger}
  \date{September 2026}
  \hypersetup{
    pdftitle={Exponential-Range Mass Production of Boolean Functions: Local recovery and optimal quantum state preparation},
    pdfauthor={Samuel Schlesinger}
  }
\fi

\begin{document}
\maketitle

\begin{abstract}
The hardest Boolean functions on $n$ bits require $\Theta(2^n/n)$ gates
to evaluate once.  Must evaluating the same function on $t$ unrelated
inputs cost $t$ times as much?  We show that every Boolean function can
be evaluated on any $t\le2^{\gamma n}$ inputs using a total of
$O_\gamma(2^n/n)$ gates, for every fixed $0\le\gamma<1$.
Thus exponentially many evaluations fit within a constant factor of the
worst-case cost of one.  In the standard AND, OR, NOT basis, our bound
is $(1/(1-\gamma)+o_\gamma(1))2^n/n$ gates.  Uhlig previously achieved
$(1+o(1))2^n/n$ gates for the smaller range $t=2^{o(n/\log n)}$.
We also attain $O_\gamma(2^n/n)$ gates and $O_\gamma(n)$ depth
simultaneously.

Our construction encodes shorter restrictions of the function and recovers
the requested outputs from disjoint sets of encoded functions.  The gate
bound includes choosing those sets, evaluating the functions, and routing
the answers.  Given the truth table, a randomized algorithm constructs
circuits attaining the simultaneous size and depth bounds in polynomial
time in the table length, with high probability.  Efficient deterministic
construction of their fixed scheduling menus remains open.  A Lean
companion verifies the recursive and sharp-coefficient Boolean size bounds.

As an application, we determine the optimal worst-case cost of preparing
$t$ copies of a specified $n$-qubit pure state throughout the same range:
$\Theta_\gamma((2^n/n)\log_2(t+1))$ elementary Clifford+$T$ gates.
The error is at most $1/10$ in trace distance for the entire batch, and
scratch qubits are restored to zero.  Coherent Boolean evaluation gives
the upper bound; a packing and circuit-counting argument gives the
matching lower bound.
We also apply the Boolean theorem to local hardness amplification:
an encoded function can remain hard to approximate while exponentially
many evaluations share the synthesis cost of its source truth table.
\end{abstract}

\section{Introduction}

How much circuit size can be saved when the same function must be evaluated on
many unrelated inputs?  Each output depends on a separate input block, but
the circuit may reorganize and share intermediate computations across copies.
This is the direct-sum question studied here; the input values may repeat.

For a Boolean function $f:\bits^n\to\bits$ and an integer $t\ge 1$, define
\[
  f^{\times t}(x^{(1)},\ldots,x^{(t)})
  =\bigl(f(x^{(1)}),\ldots,f(x^{(t)})\bigr).
\]
The naive construction uses $t$ disjoint copies of a circuit for $f$ and gives
\[
  \C(f^{\times t})\le t\C(f).
\]
Here $\C$ counts elementary Boolean gates.  The classical
Shannon--Lupanov bounds say that the hardest functions on $n$ bits have
circuit size $\Theta(2^n/n)$; in our basis the worst-case size is
$(1+o(1))2^n/n$ \cite{Shannon1949,Lupanov1958,FrandsenMiltersen2005}.
We call this the worst-case cost of one evaluation.
The mass-production problem asks how large $t$ can become before the cost must
rise above that scale.  For unrestricted circuits, the
naive direct-sum bound can be far from optimal.  Uhlig's theorem allows
$t=2^{o(n/\log n)}$ independent evaluations with no asymptotic increase in
the leading worst-case one-copy synthesis cost
\cite{Uhlig1992,Wegener1987,Kretschmer2023}.

We prove, by a different construction, that the order-of-growth plateau reaches
$2^{\gamma n}$ simultaneous evaluations for every fixed $\gamma<1$.  The
asymptotic coefficient is at most $1/(1-\gamma)$ in our basis.  This does not
preserve Uhlig's coefficient $1$ for positive fixed $\gamma$.

\subsection{Boolean size and depth bounds}

We fix the standard basis $\{\wedge,\vee,\neg\}$, with fan-in two for the
binary gates.  Constant sources, fan-out, and output designations are free;
circuit size counts the $\wedge$, $\vee$, and $\neg$ gates.
Any other fixed finite complete basis changes
the result by at most a constant-factor simulation.

\begin{theorem}[Exponential-range mass production]\label{thm:main}
For every fixed $0\le\gamma<1$, there is a constant $A_\gamma$ such that, for
every integer $n\ge1$, every Boolean function $f:\bits^n\to\bits$, and every
integer $1\le t\le 2^{\gamma n}$,
\[
  \C(f^{\times t})\le A_\gamma\frac{2^n}{n}.
\]
\end{theorem}

The next theorem refines the leading constant in this bound.

\begin{theorem}[Quantitative leading coefficient]\label{thm:coefficient}
For every fixed $0\le\gamma<1$, uniformly over Boolean functions on $n$ bits
and integers $1\le t\le2^{\gamma n}$,
\[
  \C(f^{\times t})
  \le\left(\frac{1}{1-\gamma}+o_\gamma(1)\right)\frac{2^n}{n}.
\]
Equivalently, for every $\varepsilon>0$, the coefficient
$1/(1-\gamma)+\varepsilon$ suffices for all sufficiently large $n$.
\end{theorem}

The order-of-growth size bound can also be attained at linear depth.

\begin{theorem}[Simultaneous size and depth]\label{thm:linear-depth}
For every fixed $0\le\gamma<1$, there are constants $A_\gamma,B_\gamma$
such that, for every $n\ge1$, every Boolean function $f:\bits^n\to\bits$,
and every integer $1\le t\le2^{\gamma n}$, there is a circuit $C$ computing
$f^{\times t}$ with
\[
  \operatorname{size}(C)\le A_\gamma\frac{2^n}{n},
  \qquad
  \operatorname{depth}(C)\le B_\gamma n.
\]
\end{theorem}

This simultaneous bound preserves the order of size in \cref{thm:main};
its proof does not retain the leading coefficient of \cref{thm:coefficient}.
\Cref{cor:randomized-construction} gives a randomized algorithm that,
given the full truth table, constructs such a circuit in time polynomial
in $2^n$, with probability at least $1-2^{-n}$ of correctness on every
input batch.

For fixed $\gamma$, the bound is independent of the batch size throughout
the stated range.  At the largest allowed batch size, the amortized size per
requested value is therefore
\[
  \frac{\C(f^{\times t})}{t}
  =O_\gamma\!\left(\frac{2^{(1-\gamma)n}}{n}\right).
\]
This is a worst-case synthesis guarantee.  For a particular function with a
small circuit, the replication bound $t\C(f)$ may be smaller; one may always
take the better of the two bounds.

There must nevertheless be a transition before the $2^n$-copy scale.  If
$n\ge2$ and $f$ depends on all $n$ variables, every circuit for
$f^{\times t}$ has size at least $t(n-1)$.  Hence a uniform $O(2^n/n)$ plateau
cannot persist once $t=\omega(2^n/n^2)$; see \cref{sec:discussion}.  Our
theorem reaches every fixed exponential rate, but leaves this near-endpoint
regime open.

\subsection{An optimal quantum state preparation tradeoff}

Mass production also matters when many copies of the same quantum state
are needed.  A synthesis circuit may be designed for the specified state,
and sharing can occur while the copies are prepared from all-zero inputs.
Kretschmer established quantum analogues of Uhlig's theorem by mass
producing conditional rotations and using them to prepare states
\cite{Kretschmer2023}.  Building on this framework and the connection to
coherent data loading \cite{HugginsKhattarWiebe2025}, our stronger Boolean
theorem yields an upper bound for approximate state preparation in a finite
gate set throughout the exponential copy range.  A packing lower bound
gives the matching tradeoff.

Let $M(n,t)$ be the maximum, over $n$-qubit pure states $|\psi\rangle$,
of the minimum number of elementary Clifford+$T$ gates needed to prepare
$|\psi\rangle^{\otimes t}$ to trace-distance error at most $1/10$.
The error is measured on all $nt$ output qubits together: no measurement
on the batch can change an outcome probability by more than $1/10$.
Gates may connect any qubits; scratch qubits start and end in zero.  We count every gate,
including Clifford gates, and permit circuits tailored to $\psi$.
\Cref{sec:quantum} gives the precise model and the proof.

\begin{theorem}[Optimal approximate quantum state mass production]
\label{thm:quantum}
For every fixed $0\le\gamma<1$, uniformly over integers
$1\le t\le2^{\gamma n}$ as $n\to\infty$,
\[
  M(n,t)=\Theta_\gamma\!\left(\frac{2^n}{n}\log_2(t+1)\right).
\]
\end{theorem}

The cost grows logarithmically with the number of copies.  In particular,
it is $\Theta(2^n/n)$ for one copy, $\Theta(2^n\log n/n)$ for
$t=n^a$ with fixed $a>0$, and $\Theta_\gamma(2^n)$ for
$t=\lfloor2^{\gamma n}\rfloor$ with fixed $0<\gamma<1$.
The single-copy endpoint is already known \cite{GossetKothariWu2026}.
Our upper bound uses coherent Boolean mass production and carefully
allocated rotation precision.  The lower bound reflects the fact that
$t$ copies distinguish states whose single-copy distance is on the order
of $1/\sqrt t$.

\subsection{How the construction works}

Split an input as $x=(a,z)$, with a $p$-bit prefix $a$ and a $d$-bit suffix
$z$, where $p+d=n$.  The prefix selects one of the $2^p$ shorter functions
$f_a(z)=f(a,z)$.  The desired saving would come from using only one circuit
for each shorter function.  But two inputs with the same prefix may require
that function at two different suffixes, and one circuit has only one set of
input wires.

Coding supplies alternative ways to answer each request.  We replace the
restrictions $f_a$ by a larger bank of shorter functions, called
\emph{resources}.  Each desired value can be recovered from many different
sets of resources.  A scheduler chooses one set for each request so that no
resource is needed twice.  Each resource circuit can then receive the one
suffix assigned to it, and all requested answers can be reconstructed.

Uhlig serves two requests using opposite ends of an XOR chain, then scales
through recursion \cite{Uhlig1992,Wegener1987}.  We replace that chain with a
code offering many recovery choices, and schedule disjoint choices for an
entire batch.  The key is making the scheduling cheap enough to preserve the
savings.  \Cref{sec:two-copy} develops this comparison through a worked example.

The construction has four components.
\begin{enumerate}[label=\textbf{\arabic*.},leftmargin=2.2em]
\item \textbf{Pack restrictions.}
For every fixed suffix $z$, encode the values $f(a,z)$ over all prefixes $a$
in a code over $\F_q$.  As $z$ varies, each bit of each code coordinate
defines a resource function on $d$ variables.

\item \textbf{Recover locally.}
The code coordinates are points in $\F_q^\ell$.  A requested symbol is the
sum of the $q-1$ other symbols on any line through its information point.
About $q^{\ell-1}$ directions are available, giving many recovery options.

\item \textbf{Schedule disjoint recovery.}
Test a fixed menu of direction assignments and choose one that serves at
least half the pending requests without collisions.  Reserve those recovery
sets and repeat on the remaining requests.  The total gate cost is
$\widetilde O_\ell(tq)$, including selection of a schedule for the live inputs.

\item \textbf{Evaluate the resources.}
Each code coordinate is now queried on at most one suffix.  Apply ordinary
one-copy synthesis to each resource function on $d=|z|$ variables, route its
value to the requesting line, and decode.
\end{enumerate}

The encoding defines the resource truth tables when the circuit is built.
At evaluation time, the circuit selects recovery sets, routes suffixes,
evaluates the fixed resource circuits, and decodes.  The menus are fixed in
advance and work for every batch; the input prefixes determine which entries
the circuit selects.  The depth refinement in \cref{sec:linear-depth}
increases the progress per phase so that only constantly many phases remain.

\subsection{Where the leading coefficient comes from}

The main cost is the product of the number of resource functions and the cost
of synthesizing one of them.  A code whose rate tends to one, together with
packing that leaves only a vanishing fraction of its capacity unused, gives
\[
  \underbrace{(1+o(1))2^p}_{\text{Boolean resources}}
  \;\cdot\;
  \underbrace{(1+o(1))\frac{2^d}{d}}_{\text{gates per resource}}
  =(1+o(1))\frac{2^n}{d}.
\]
Leaving more suffix bits makes this product smaller, but leaves fewer prefix
bits to provide distinct recovery options.  The scheduler can accommodate
$t\le2^{\gamma n}$ while leaving $d=\delta n+O(1)$ suffix bits for any fixed
$0<\delta<1-\gamma$, by taking the geometric dimension $\ell$ sufficiently
large and then fixed.  Its scheduling and routing costs are exponentially
smaller than $2^n/n$.  The displayed product therefore gives the coefficient
$1/\delta$.  Choosing $\delta$ close enough to $1-\gamma$ gives
$1/(1-\gamma)+\varepsilon$ for any fixed $\varepsilon>0$.
For example, the bound for $t\le2^{n/2}$ has coefficient $2+o(1)$.
\Cref{app:module-bound} explains when this coefficient is forced by the
resource-bank architecture.

\paragraph{Machine-checked companion.}
A Lean~4 companion verifies the recursive size bound and the complete
nonuniform construction yielding \cref{thm:coefficient}, including its
real-rate and additive-error quantifiers.  \Cref{sec:formal-accounting}
records the theorem endpoints, pinned sources, and finite accounting.
The depth, randomized-construction, quantum, local-encoding, and lower-bound results
are written proofs outside this companion.

\paragraph{Development methodology.}
Generative-AI systems materially assisted mathematical experiments,
literature discovery, and proof development, as detailed in the AI Disclosure
at the end of the paper.

\paragraph{Organization.}
\Cref{sec:related,sec:conventions} compare prior work and fix the circuit
model.  \Cref{sec:two-copy} illustrates disjoint recovery with two requests.
\Cref{sec:code,sec:scheduling,sec:master} develop the code, scheduler, and
composition bound for the direct proof in \cref{sec:direct}.
\Cref{sec:high-rate,sec:linear-depth} refine the coefficient and depth;
the latter also gives the randomized construction.
\Cref{sec:quantum} proves the quantum state preparation tradeoff.
\Cref{sec:amplification} combines local hardness amplification with mass
production and relates the construction to Hirahara's partial-MCSP
reduction.  \Cref{sec:discussion} records a secure-evaluation consequence
and discusses the leading coefficient and open problems.

The appendices give an explicit recursive alternative
(\cref{sec:induction}), circuit implementation details
(\cref{app:implementation}), exact accounting for the Lean companion
(\cref{sec:formal-accounting}), and the lower bound for modules with
polynomial interfaces (\cref{app:module-bound}).

\section{Relation to earlier work}\label{sec:related}

The table compares the circuit bounds and the main coding antecedents.
Here ``one-copy leading term'' means $(1+o(1))2^n/n$ gates in our basis,
the worst-case cost of evaluating one function once.
Uhlig's 1974 paper \cite{Uhlig1974} is a precursor to
the mass-production theorem presented in \cite{Uhlig1992} and in Wegener's
Theorem~10.2.3 \cite{Wegener1987}.

\begin{center}
\small
\begin{tabularx}{\linewidth}{@{}
  >{\raggedright\arraybackslash}p{0.14\linewidth}
  >{\raggedright\arraybackslash}p{0.20\linewidth}
  >{\raggedright\arraybackslash}p{0.27\linewidth}
  >{\raggedright\arraybackslash}X@{}}
\toprule
Result & Copy range & Circuit-size bound & Comparison \\
\midrule
Uhlig \cite{Uhlig1992,Wegener1987,Kretschmer2023}
  & $\log t=o(n/\log n)$
  & one-copy leading term
  & optimal order and leading constant in a smaller copy range \\
Paul, Lemma~2 \cite{Paul1976}
  & arbitrary $t$
  & $O(\max\{n2^n,nt\log^2t\})$
  & $O(n2^n)$ throughout every fixed exponential range below one \\
Holmgren--Rothblum \cite{HolmgrenRothblum2024}
  & arbitrary $t$
  & $O(2^n+tn^3)$ after hardwiring
  & $O(2^n)$ throughout every fixed exponential range below one \\
Albrecht, Theorem~1 \cite{Albrecht1989}
  & $t\le2^n$
  & $O(2^n n^2)$
  & the full exponential range, with polynomial overhead \\
Finite-geometry batch codes \cite{IKOS2004,PolyanskiiVorobyev2019}
  & code-dependent batch size
  & disjoint local recovery; no bound on $\C(f^{\times t})$
  & coding antecedents for affine-geometric recovery \\
This paper
  & $t\le2^{\gamma n}$ for fixed $\gamma<1$
  & $O_\gamma(2^n/n)$
  & optimal order; limiting coefficient at most $1/(1-\gamma)$ \\
\bottomrule
\end{tabularx}
\end{center}

Paul's construction treats the truth table as an ordered database: sort the
input addresses, sweep through the hardwired table, and restore the answers
to their original order.  For fixed $\gamma<1$ and $t\le2^{\gamma n}$,
$nt\log^2t=o(n2^n)$, giving the table's $O(n2^n)$ bound.
Paul's $g^k$ is the same product task considered here; his arbitrary binary
Boolean gates can be simulated in our basis at constant cost.

Holmgren and Rothblum prove that $Q$ indexed bits can be selected from a live
$N$-bit database by a bounded-fan-in Boolean circuit of size
$O(N+Q\log^3N)$ and depth $O(\log(N+Q))$
\cite[Main theorem]{HolmgrenRothblum2024}.  In our nonuniform model, hardwiring
the database to the $2^n$-bit truth table of $f$ and using the $t$ live
$n$-bit input blocks as indices gives $O(2^n+tn^3)$, as recorded in the table.
Thus, throughout every fixed exponential range below one, \cref{thm:main}
improves the bounds from Paul's construction and this multiselection reduction
by factors of $n^2$ and $n$, respectively.  Multiselection also implements the
gather step of our routing; \cref{app:routing-proof} gives the reduction.

Albrecht studies the same simultaneous-realization problem over a complete
binary basis \cite[Theorem~1, pp.~53--54]{Albrecht1989}.  His construction
covers the endpoint $t=2^n$, with depth $O(n\log n)$ and size a factor
$O(n^3)$ above the one-copy scale, as Albrecht also notes.

Hiltgen and Paterson use ``mass production'' for $\mathrm{PI}_k$-set circuits
underlying monotone slice functions \cite{HiltgenPaterson1995}.  Their result
gives sharing in that restricted set-circuit setting and leads to an optimal
circuit for the Ne\v{c}iporuk slice.  The task here is simultaneous evaluation
of an arbitrary Boolean function on disjoint input blocks.

\paragraph{Direct sums in oracle models.}
Suruga \cite[Theorem~1]{Suruga2024} characterizes limiting amortized
worst-case oracle complexity by optimal expected one-instance oracle
complexity, for a fixed problem and positive per-instance error.
His framework charges oracle calls and leaves internal computation
uncharged.  Our circuits compute every output exactly, and their savings
concern the gates used for that internal computation, including schedule
selection and routing.

\paragraph{Coding ingredients and circuit costs.}
Systematic polynomial evaluation, affine-line recovery, and disjoint recovery
sets are established batch-code tools.  Section~3.4 of
Ishai--Kushilevitz--Ostrovsky--Sahai \cite{IKOS2004} chooses random affine lines
and deletes
intersections while preserving enough points for interpolation.  Their
Remark~3.11 also observes that the randomized decoding algorithms in that and
the following sections can be derandomized using limited independence.
In their terminology, the repeated-request analogue relevant here is a
primitive multiset batch code.  Our local gadget uses the elementary
tensor-product evaluation subspace, with a separate degree bound in each
variable, whereas their construction uses total-degree Reed--Muller evaluation.

Finite-geometry constructions of primitive batch codes likewise give explicit
families of mutually disjoint recovery sets \cite{PolyanskiiVorobyev2019}.
Later lifted-polynomial constructions also use line-based recovery to obtain
batch-code guarantees \cite{HolzbaurEtAl2020}.  Private-information-retrieval
codes of Fazeli, Vardy, and Yaakobi form a related family in which each
information symbol has many disjoint recovery options
\cite{FazeliVardyYaakobi2015}.

For circuit mass production, selecting recovery sets is itself part of the
computation: the prefixes are input wires, and every batch must produce valid
recovery records within the gate budget.  A generic polynomial bound in the
exponentially large code length and batch size need not fit below $2^n/n$.
The menu scheduler provides the more precise bound
$\widetilde O_\ell(tq)$, which \cref{prop:scheduler-sensitive} combines with
resource synthesis and routing.  Counting occupied-line descriptions keeps the
menus short; sharing the occupied-point list across candidates keeps their
verification small.  Together with high-rate packing, this gives the stated
exponential-range bound and its leading coefficient.

\paragraph{Quantum mass production and synthesis.}
Kretschmer \cite[Sections~3--4]{Kretschmer2023} gives exact synthesis of
$t=2^{o(n/\log n)}$ copies of any $n$-qubit state with $O(2^n)$ CNOT
gates, allowing arbitrary one-qubit gates.  He also obtains an $O(4^n)$
bound for $t$ uses of an arbitrary $n$-qubit unitary.  His reduction from
mass production of conditional rotations to state preparation provides
the framework for our upper bound.  His discussion asks what happens for
larger copy ranges.  \Cref{thm:quantum} determines the total elementary
gate count for $t\le2^{\gamma n}$, for every fixed $\gamma<1$, in a finite
gate set at constant error for the whole batch.  The larger copy range
and matching finite-gate bounds are the advance over his result; the
model differences remain essential to the comparison with exact synthesis.

Huggins, Khattar, and Wiebe \cite{HugginsKhattarWiebe2025} combine mass
production with quantum read-only memory, which loads fixed classical
data at addresses held in quantum registers.  They give applications to
quantum algorithms and analyze both a small-copy plateau and a
larger-copy regime.  Their Appendix~B.5 explains the passage from data
loading to state preparation.  Our upper proof uses this established
connection with the stronger Boolean synthesis bound.

Gosset, Kothari, and Wu \cite{GossetKothariWu2026} determine the optimal
$T$-gate count for state preparation and diagonal-unitary synthesis.
Their discussion also records an $O(2^n/n)$ total-gate bound for preparing
one state at constant error, and their Appendix~B allocates approximation
error geometrically across conditional rotations.  We count all elementary
gates and determine their dependence on $t$ through an exponential copy
range.  Our proof uses standard conditional-rotation synthesis
\cite{ShendeBullockMarkov2006} and finite-gate approximation
\cite{DawsonNielsen2006}, followed by a matching packing lower bound.

\section{Circuit conventions and worst-case notation}\label{sec:conventions}

Our circuits are finite directed acyclic graphs.  Source vertices are input
bits or the constants zero and one.  Internal vertices are labelled by
$\wedge$, $\vee$, or $\neg$, with the usual fan-in.  A circuit for a map
$F:\bits^N\to\bits^m$ has $m$ designated output wires, which may be input
wires or gate outputs and may repeat.  Fan-out, constants, and output
designations carry no cost; size is the number of internal gates.  We write
$\C(F)$ for the minimum size of such a circuit.  All leading-constant
comparisons in this paper refer to this fixed model.  The depth of a circuit
is the maximum number of internal gates on a path from a source to a
designated output, counting a $\neg$ gate as one gate on such a path.

Let
\[
  L(n)=\max_{f:\bits^n\to\bits}\C(f).
\]
The Shannon counting lower bound and Lupanov's synthesis upper bound imply
\[
  L(n)=\Theta(2^n/n)
\]
for every fixed finite complete basis; see
\cite{Shannon1949,Lupanov1958,FrandsenMiltersen2005,Wegener1987}.

Fixing all but one input block and retaining its output gives
$\C(f^{\times t})\ge\C(f)$.  Thus $2^n/n$ is the optimal worst-case
order even when several outputs are requested.

Circuits may depend on the full truth table of $f$.  Their size and depth
measure evaluation after the circuit has been constructed.  Field
representations, irreducible polynomials, sorting networks, and scheduling
menus may all be fixed as part of that construction.
\Cref{cor:randomized-construction} separately bounds the time to construct
the circuits for the simultaneous size and depth theorem from a supplied
truth table.

The parameters of the Boolean construction have the following roles.
Code length and dimension are measured in field symbols; multiplying by
$b$ converts either count to bits.
\begin{center}
\small
\begin{tabularx}{\linewidth}{@{}lX@{}}
\toprule
Parameter & Meaning \\
\midrule
$n=p+d$, $t$ & Input length, prefix/suffix split, and number of requests \\
$q=2^b$, $\ell$ & Field size, bits per symbol, and fixed geometric dimension \\
$N=q^\ell$, $K$ & Code length and dimension; $K/N$ is the code rate \\
$D_q=(q^\ell-1)/(q-1)$ & Number of line directions through each target \\
$g$ & Number of requests in the scheduler lemmas \\
$M$ & Number of slots in a routing table \\
\bottomrule
\end{tabularx}
\end{center}

We use $\widetilde O_\ell(X)$ for
\[
  X\cdot\poly_\ell(n,\log q,\log(t+2),\log(M+2)).
\]
Only polynomial factors are suppressed; in the parameter regimes used below,
they never affect an exponential rate.

For a positive integer $k$, write $[k]_0=\{0,\ldots,k-1\}$.

We repeatedly use two standard circuit primitives.  First, after representing
$\F_{2^b}$ in a polynomial basis, field addition, multiplication, inversion of
a nonzero element, and comparison of canonical $b$-bit representations have
circuits of size $\poly(b)$.  For example, inversion can be implemented by
binary exponentiation to the power $2^b-2$.  Second, $R$ records of width $w$
can be sorted by a fixed sorting network using $O(R\log^2 R)$ compare--exchange
units, each of size $O(w)$; see \cite{Batcher1968}.  Sorting by a leading type
bit, with original positions used as tie breakers when needed, also gives
stable fixed-wire compaction within the same asymptotic bound.

\section{The two-copy construction as disjoint recovery}\label{sec:two-copy}

Uhlig's two-copy argument illustrates the invariant behind disjoint recovery:
each requested value has two representations, and the requests can choose
representations that use separate resources.

Fix a prefix length $p$, put $M=2^p$, and define the restrictions
\[
  f_i(z)=f(i,z),
  \qquad i\in\{0,\ldots,M-1\}.
\]
Uhlig's two-copy construction introduces the following functions; see the proof
of Wegener's Theorem~10.2.3 \cite{Wegener1987} and the overview in
\cite[Section~1.2]{Kretschmer2023}:
\[
  g_0=f_0,
  \qquad
  g_i=f_{i-1}\xor f_i\quad(1\le i<M),
  \qquad
  g_M=f_{M-1}.
\]
Every restriction has two representations:
\[
  f_i=\bigxor_{j=0}^{i}g_j
  =\bigxor_{j=i+1}^{M}g_j.
\]
Given two inputs $(i,z^{(1)})$ and $(j,z^{(2)})$ with $i\le j$, recover
$f_i(z^{(1)})$ from the prefix set
$\{0,\ldots,i\}$ and $f_j(z^{(2)})$ from the suffix set
$\{j+1,\ldots,M\}$.  The sets are disjoint, so each resource function $g_j$
is evaluated on at most one suffix.

For example, take $p=2$, so there are four restrictions and five resources.
The requests with prefixes $1$ and $2$ can be answered as
\[
  f_1(z^{(1)})=g_0(z^{(1)})\xor g_1(z^{(1)}),
  \qquad
  f_2(z^{(2)})=g_3(z^{(2)})\xor g_4(z^{(2)}).
\]
Resources $g_0,g_1$ receive $z^{(1)}$, resources $g_3,g_4$ receive $z^{(2)}$,
and $g_2$ is unused.  The suffixes can be completely unrelated.  Even if the
two prefixes coincide, the prefix and suffix recovery sets remain disjoint.

For two requests arriving in the opposite order, a comparator first orders
their prefixes, conditional swaps route the two suffixes to the prefix and
suffix recovery branches, and a final swap restores the output order.  This
costs only $\poly(n)$ gates.

The map $(f_0,\ldots,f_{M-1})\mapsto(g_0,\ldots,g_M)$ is, up to an invertible
change of information coordinates, the length-$(M+1)$ single-parity-check
code.  The two representations above are disjoint recovery sets for each
information coordinate.  To serve exponentially many requests, we use a
constant-rate code with many geometric recovery sets.

\section{The local-recovery gadget}\label{sec:code}

We need a code with three features: it stores the $2^p$ prefix values using
only $O(2^p)$ bits, it offers many recovery sets for every value, and the
circuit can compute the addresses in those sets cheaply.  The following
product code supplies all three.  \Cref{prop:local-gadget} is the interface
used by the rest of the proof; the polynomial construction establishes that
interface.

Fix an integer $\ell\ge2$ and a field $\F_q$ of characteristic two, where
$q=2^b$.  Set
\[
  s_0=\left\lfloor\frac{q-1}{\ell}\right\rfloor
\]
and fix a polynomial-basis representation of $\F_q$.  If $\alpha_i$ denotes
the field element whose canonical $b$-bit representation is the integer $i$,
choose the efficiently indexed set
\[
  A=\{\alpha_0,\ldots,\alpha_{s_0-1}\}.
\]
\begin{samepage}
Let $\mathcal V$ be the space of polynomials
\[
  P(X_1,\ldots,X_\ell)\in\F_q[X_1,\ldots,X_\ell]
\]
having degree less than $s_0$ in every variable.  Repeated univariate
interpolation shows that evaluation on $A^\ell$ uniquely determines $P$, and
\[
  \deg P\le\ell(s_0-1)\le q-2.
\]
\end{samepage}

The evaluation map
\[
  \operatorname{Enc}:\F_q^{A^\ell}\longrightarrow\F_q^{\F_q^\ell},
  \qquad
  \bigl(P(u)\bigr)_{u\in A^\ell}
  \longmapsto
  \bigl(P(y)\bigr)_{y\in\F_q^\ell},
\]
is a systematic linear code: the original information symbols appear
unchanged at the coordinates in $A^\ell$.  Its length and dimension are
\[
  N=q^\ell,\qquad K=s_0^\ell=\Theta_\ell(q^\ell).
\]

\subsection{Packing Boolean restrictions}\label{subsec:packing}

Split an $n$-bit input as
\[
  x=(a,z),
  \qquad |a|=p,
  \qquad |z|=d=n-p.
\]
For each fixed suffix $z$, pack the $2^p$ bits
\[
  \bigl(f(a,z)\bigr)_{a\in\bits^p}
\]
into $K$ field symbols, using the fixed $\F_2$-basis of $\F_q$ and padding
unused bit positions with zeros.  More explicitly, interpret a prefix $a$ as
an integer $I(a)\in[2^p]_0$, and set
\[
  s(a)=\left\lfloor\frac{I(a)}b\right\rfloor,
  \qquad
  j(a)=I(a)\bmod b.
\]
Write $s(a)$ in base $s_0$ using exactly $\ell$ digits
$i_1(a),\ldots,i_\ell(a)\in[s_0]_0$, and define
\[
  u(a)=\bigl(\alpha_{i_1(a)},\ldots,\alpha_{i_\ell(a)}\bigr)\in A^\ell.
\]
Thus bit $j(a)$ of the information symbol at $u(a)$ is $f(a,z)$.
Fixed-divisor division, remainder, and base conversion have Boolean circuits
of size $\poly_\ell(p,\log q)$, so computing $(u(a),j(a))$ for all requests
costs $\widetilde O_\ell(t)$ gates.

Choose $q$, among powers of two, minimally so that
\[
  Kb\ge2^p.
\]
For all sufficiently large $p$ (depending on $\ell$), the corresponding code
dimensions at field sizes $q$ and $q/2$ are both $\Theta_\ell(q^\ell)$.
Minimality then gives
\begin{equation}\label{eq:packing}
  q^\ell b=\Theta_\ell(2^p).
\end{equation}
Indeed, the lower bound follows from $Kb\le q^\ell b$.  For the upper bound,
the preceding candidate $q/2$ has bit width $b-1$, fails the packing
inequality, and has dimension $\Omega_\ell(q^\ell)$.  Hence
$q^\ell(b-1)=O_\ell(2^p)$; since $b\le2(b-1)$ for $b\ge2$, the upper bound in
\eqref{eq:packing} follows.  Equivalently,
\begin{equation}\label{eq:q-log}
  \ell\log_2 q+\log_2\log_2 q=p+O_\ell(1),
  \qquad
  \log_2q=\frac p\ell+O_\ell(\log p).
\end{equation}
Let $P_z$ be the polynomial whose information symbols are the packed bits.
For each code coordinate $y\in\F_q^\ell$, define
the field-valued resource
\[
  G_y(z)=P_z(y).
\]
Each of the $b$ coordinate bits of $G_y$ is a Boolean function on $d$
variables, determined by $f$, $y$, and the chosen field basis.
For $r\in[b]_0$, write
\[
  H_{y,r}(z)=\bigl(G_y(z)\bigr)_r.
\]

These definitions view the same array in two ways:
\begin{center}
\begin{tabular}{@{}lll@{}}
\toprule
Held fixed & What varies & What the values describe \\
\midrule
Suffix $z$ & Code coordinate $y$ & One codeword $P_z$ \\
Code coordinate $y$ & Suffix $z$ & One resource function $G_y$ \\
\bottomrule
\end{tabular}
\end{center}
The decoder combines coordinates of a codeword.  The resource circuit
evaluates one column of this array at the suffix routed to it.  No encoding
circuit or full array is evaluated on the live inputs: the array specifies
which shorter functions are synthesized.

\subsection{Affine-line recovery}

\begin{lemma}[Line identity]\label{lem:line}
Let $P:\F_q^\ell\to\F_q$ be represented by a polynomial of total degree at
most $q-2$.  For every $u\in\F_q^\ell$ and every nonzero direction
$v\in\F_q^\ell$,
\[
  P(u)=-\sum_{\lambda\in\F_q^*}P(u+\lambda v).
\]
In characteristic two, the minus sign may be omitted.
\end{lemma}

\begin{proof}
The restriction $Q(\lambda)=P(u+\lambda v)$ is a univariate polynomial of
degree at most $q-2$.  The sum of the constant monomial over $\F_q$ is
$q\cdot1=0$ in characteristic two.  For $1\le j\le q-2$, the sum of
$\lambda^j$ over $\F_q^*$ vanishes because $j$ is not divisible by $q-1$,
and the zero term contributes nothing.  Applying these identities monomial by
monomial gives
$\sum_{\lambda\in\F_q}Q(\lambda)=0$.  Separating the term $\lambda=0$ proves
the claim.
\end{proof}

Thus every affine line through an information point $u$ supplies a recovery
set consisting of its $q-1$ non-target points.  We identify two nonzero
direction vectors when one is a nonzero scalar multiple of the other and
write $[v]$ for the resulting projective direction.  The number of such
directions through a point is
\begin{equation}\label{eq:directions}
  D_q=\frac{q^\ell-1}{q-1}=1+q+\cdots+q^{\ell-1},
  \qquad
  q^{\ell-1}\le D_q<2q^{\ell-1}.
\end{equation}
For fixed $\ell$ and $p\to\infty$, \cref{eq:q-log,eq:directions} give the
parameter dictionary
\begin{equation}\label{eq:parameter-dictionary}
  q^\ell=2^{p+o(p)},
  \qquad
  q=2^{p/\ell+o(p)},
  \qquad
  D_q=q^{\ell-1+o(1)}
      =2^{((\ell-1)/\ell)p+o(p)}.
\end{equation}
Two distinct lines through the same target intersect only at that target,
which is absent from both recovery sets.  Repeated requests for one information
symbol therefore cause no additional difficulty.  A later target may itself
belong to the occupied set, but it blocks no direction because every new
recovery set omits its own target.

The facts needed later can now be packaged without reference to the internal
polynomial representation.

\begin{proposition}[Local-recovery gadget]\label{prop:local-gadget}
Fix $\ell\ge2$ and a split $n=p+d$, with $p$ sufficiently large depending on
$\ell$, and choose $q=2^b$ by the packing rule above.  For every
$f:\bits^n\to\bits$ there are $q^\ell b$ Boolean resource functions
$H_{y,r}:\bits^d\to\bits$ with the following properties.
\begin{enumerate}[label=\textup{(\roman*)},leftmargin=2.1em]
\item The resource count satisfies $q^\ell b=\Theta_\ell(2^p)$.
\item A prefix $a$ determines an information point $u(a)\in A^\ell$ and a bit
position $j(a)\in[b]_0$ by a circuit of size $\poly_\ell(p,\log q)$.
\item For every projective direction $[v]$, the set
\[
  R(u(a),[v])=\{u(a)+\lambda v:\lambda\in\F_q^*\}
\]
has $q-1$ coordinates and recovers the requested bit by
\[
  f(a,z)=\bigxor_{y\in R(u(a),[v])}H_{y,j(a)}(z).
\]
There are $D_q=(q^\ell-1)/(q-1)$ such recovery sets.
\item If $U\subseteq\F_q^\ell$ is already occupied, at most
$|U\setminus\{u(a)\}|$ directions have
$R(u(a),[v])\cap U\ne\varnothing$.
\end{enumerate}
\end{proposition}

\begin{proof}
The packing construction gives (i) and (ii).  The line identity, followed by
selection of coordinate bit $j(a)$ in the fixed $\F_2$ basis, gives (iii).
For (iv), every $y\in U\setminus\{u(a)\}$ lies on exactly one projective line
through $u(a)$, while the omitted target $u(a)$ blocks none.
\end{proof}

Thus one desired bit has become $q-1$ requests to shorter resource functions,
but those coordinates may be chosen from any of $D_q$ recovery sets.  The next
section turns the blocking guarantee in (iv) into a deterministic circuit and
then routes only the coordinates actually selected.

\section{Deterministic scheduling and routing}\label{sec:scheduling}

The scheduler maps target coordinates to disjoint recovery sets by testing
fixed menus, using $\widetilde O_\ell(gq)$ gates for $g$ requests.  Routing
then connects the selected addresses to the resource circuits.

The blocking bound also supports a greedy approach: occupied points rule out
directions, and a circuit selects the first available direction for each
request.  Its $\widetilde O_\ell(g^2q)$ cost leads to the recursive
construction in \cref{sec:induction}, where this scheduler is developed.

\subsection{A nonuniform scheduler with linear dependence on the batch size}
\label{subsec:linear-scheduler}

The menu scheduler works in phases that accept half the remaining requests.
We prove that a small fixed menu of candidate directions works for
every possible phase state, then give a circuit that tests the menu.  The
existence proof uses randomness only to choose constants of the circuit; its
evaluation is deterministic on every input.

There are three steps.  First, one random direction assignment succeeds with
very high probability for any fixed state.  Second, a union bound supplies a
short menu that contains a successful assignment for every state.  Third, a
sorting circuit finds such an assignment without multiplying the occupied-set
cost by the menu length.  Here $g$ is the original batch size, $k$ is the
number of pending requests, and $U$ contains the points reserved earlier.
A menu entry is a tuple of $k$ directions, one for each pending request.

\begin{lemma}[Collision tail]\label{lem:collision-tail}
Let $1\le k\le g$, let $U$ be a union of at most $g$ recovery sets, and
fix any ordered list of $k$ target points.  Choose their directions
independently and uniformly.  A request is bad if its candidate recovery
set meets $U$ or another candidate recovery set.
For every integer $0\le v<k$, put $h=\lceil(v+1)/2\rceil$.  Then
\[
  \Pr\{\text{more than }v\text{ requests are bad}\}
  \le \binom{k}{h}\left(\frac{2gq}{D_q}\right)^h.
\]
\end{lemma}

\begin{proof}
For any fixed set $B$, a random recovery set through a fixed target meets
$B$ with probability at most $|B|/D_q$: each point other than the target
blocks exactly one direction.  Pairwise collision events can be dependent.
We handle this by selecting a set of requests with witnesses outside it
and conditioning on the other directions.

Form a graph on the requests and a vertex $0$ representing $U$.  Join a
request to $0$ when its recovery set meets $U$, and join two requests when
their recovery sets meet.  The bad requests are precisely the nonisolated
vertices other than $0$.  If more than $v$ requests are bad, two-color a
spanning forest of the nontrivial components.  One color contains at least
$h$ bad requests, not counting vertex $0$.  Choose exactly $h$ of them as
$S$.  Every member of $S$ has a forest neighbor outside $S$, which may be
vertex $0$.

For a fixed $S$, condition on all directions outside $S$.  Their recovery
sets together with $U$ contain at most $2g(q-1)$ points.  Each target in
$S$ independently chooses a direction whose recovery set meets this fixed
union with probability at most $2gq/D_q$.  Averaging over the outside
directions and taking a union bound over the $\binom{k}{h}$ choices of
$S$ gives the estimate.
\end{proof}

\begin{lemma}[A partial schedule with exponentially small failure]
\label{lem:partial-schedule}
Put $N=q^\ell$.  Suppose $512gq\le D_q$, let $1\le k\le g$, and let $U$ be
the union of at most $g$ affine-line recovery sets.  Fix any $k$ targets, and
choose one projective direction independently and uniformly for each target.
Call a candidate recovery set \emph{clean} if it avoids $U$ and every other
candidate recovery set.  Then
\[
  \Pr\{\text{fewer than }\lceil k/2\rceil\text{ sets are clean}\}
  \le 2^{-k}.
\]
\end{lemma}

\begin{proof}
Take $v=\lfloor k/2\rfloor$ in \cref{lem:collision-tail}.  Fewer than
$\lceil k/2\rceil$ clean sets means more than $v$ bad requests.  Since
$h=\lceil(v+1)/2\rceil\ge k/4$ and $2gq/D_q\le1/256$, the failure
probability is at most
\[
 \binom{k}{h}\left(\frac{2gq}{D_q}\right)^h
 \le2^k(1/256)^{k/4}=2^{-k}.
\]
\end{proof}

\begin{lemma}[Universal menus]\label{lem:universal-menus}
Under the same slack condition, for each $1\le k\le g$ there is a fixed menu
of
\begin{equation}\label{eq:menu-size}
  r_k=\left\lceil\frac{g(1+3\log_2N)+1}{k}\right\rceil
\end{equation}
direction $k$-tuples such that every set $U$ described in
\cref{lem:partial-schedule} and every ordered list of $k$ targets has a menu
entry with at least $\lceil k/2\rceil$ clean recovery sets.
\end{lemma}

\begin{proof}
A recovery set is specified by an anchor and a direction, giving at most
$ND_q$ descriptions.  Allowing an unused marker in each of $g$ slots, the
number of occupied-set descriptions is at most $(1+ND_q)^g$.  The targets
have at most $N^k$ descriptions.  Since $D_q\le N$ and $k\le g$, the total
number of states is at most
\[
  (1+ND_q)^gN^k\le2^{g(1+3\log_2N)}.
\]
Counting line descriptions is essential: treating $U$ as an arbitrary list
of $g(q-1)$ points would introduce an extra factor of $q$ in this exponent.
Choose $r_k$ candidate tuples independently.  For a fixed state, the
probability that every tuple fails is at most $2^{-kr_k}$ by
\cref{lem:partial-schedule}.  A union bound over states is strictly below
one by \eqref{eq:menu-size}.  Fix a menu working for all states as part of
the circuit for parameters $g,k,q,\ell$.
\end{proof}

\begin{samepage}
\begin{theorem}[Nonuniform linear scheduling]\label{thm:linear-scheduler}
If $512gq\le D_q$, a deterministic nonuniform Boolean circuit maps every
ordered multiset of $g$ targets in $\F_q^\ell$ to enumerated pairwise disjoint
affine-line recovery sets using $\widetilde O_\ell(gq)$ gates.
More explicitly, its size is $O(gq\Lambda^5)$ for
$\Lambda=\ell\log_2q+\log_2(g+2)$.
The circuit's fixed menus have total description length
$O(g(\log N)^2\log(g+2))$ bits.
\end{theorem}
\end{samepage}

\begin{proof}
Maintain $k$ active targets and the occupied set from the $g-k$ accepted
requests.  Initially $k=g$ and $U$ is empty.  Test the menu from
\cref{lem:universal-menus}, choose its first successful entry, and accept
exactly the first $\lceil k/2\rceil$ clean requests in that entry.  These
sets avoid $U$ and each other.  Compact the other $\lfloor k/2\rfloor$
requests, retaining their original identifiers.  This gives a fixed sequence
of at most $1+\lfloor\log_2 g\rfloor$ phase sizes, and never occupies more
than $g$ recovery sets.  After the last phase, sort accepted records by original
identifier to enumerate the answer sets in request order.

It remains to bound the cost of testing a menu.  Generate records for all
$r_k k(q-1)$ candidate points, labelled by candidate, request, and line
position, and append the $(g-k)(q-1)$ occupied points \emph{once}.  Since
$r_k k=O(g\log N)$, there are $O(gq\log N)$ records.  Use these fixed passes:
\begin{enumerate}[label=\textup{(\roman*)},leftmargin=2.1em]
\item Sort by point, with occupied records first at equal points.  Along each
equal-point run propagate an occupancy bit.  In sorted position $i$, the
update is $e_i=\mathrm{occupied}_i\vee
((y_i=y_{i-1})\wedge e_{i-1})$, with a dummy predecessor at the boundary.
Thus every candidate point learns whether it lies in $U$.
\item Compact candidate records and sort them by $(\text{candidate},y)$.
Mark a record when its predecessor or successor has the same key.  Points on
one recovery line are distinct, so such a match is exactly an intersection
between two requests in the same candidate.
\item Sort by $(\text{candidate},\text{request},\text{line position})$ and
OR the $q-1$ bad-point flags for each request.  Count the clean requests in
each candidate and use priority selection to choose its first successful
entry.  Fixed multiplexers and compaction implement the acceptance step.
\end{enumerate}
Every record has polynomial width in $\ell\log q$ and $\log(g+2)$, and a
constant number of sorting networks and linear scans suffices.  Thus a phase
has size $O(gq\Lambda^4)$: there are $O(gq\Lambda)$ records, with
$O(\Lambda)$ bits per point record and $O(\Lambda^2)$ sorting-network
comparators per record.  Candidate directions and line scalars are hardwired,
so generating a candidate point adds a constant field vector to the target
using $O(\ell\log q)$ gates.  Selecting and compacting whole line records
costs at most $O(gq\Lambda^3)$ additional gates.  There are $O(\Lambda)$
phases, proving the explicit bound.  The menus contain
$O(g\log N)$ directions per phase,
each of $O(\log N)$ bits, giving the description bound.
\end{proof}

\subsection{Scatter and gather}

The scheduler produces a sparse list of resource addresses, whereas the
resource circuits occupy fixed positions in the circuit.  \emph{Scatter}
sends each requested suffix to its resource position; \emph{gather} returns
the resulting value to the request that needs it.  Sorting implements both
maps with cost proportional, up to polynomial factors, to the lists being
connected.  With globally disjoint recovery sets, each resource position
receives at most one suffix.

Every non-target point of a selected line produces an incidence record
containing a request identifier, a line position, the point's code
coordinate $y$, and the suffix to be evaluated.  There are exactly
$t(q-1)<tq$ such records, with distinct code coordinates.

The routing consists of four fixed sorting-and-local-update passes.  Here
``key'' is the resource coordinate $y$; any unlisted tie is broken by the record's
original position.  The record-type order is pass-specific and is shown
explicitly in the table.
\begin{center}
\small
\begin{tabularx}{\linewidth}{@{}
  >{\raggedright\arraybackslash}p{0.14\linewidth}
  >{\raggedright\arraybackslash}p{0.21\linewidth}
  >{\raggedright\arraybackslash}p{0.27\linewidth}
  >{\raggedright\arraybackslash}X@{}}
\toprule
Pass & Records & Sort order & Local result \\
\midrule
Scatter match
  & incidence, slot
  & $(\text{key},\text{type})$; incidence before slot
  & a matched slot takes its predecessor's suffix; otherwise it takes a dummy \\
Scatter order
  & incidence, filled slot
  & $(\text{type},\text{key})$; slots first
  & slots are in canonical key order \\
Gather match
  & value, saved incidence
  & $(\text{key},\text{type})$; value before incidence
  & a matched incidence takes its predecessor's value \\
Gather order
  & value, updated incidence
  & $(\text{type},\text{request},\text{line position})$; incidences first
  & decoder inputs are in canonical request order \\
\bottomrule
\end{tabularx}
\end{center}

Scatter records carry a $d$-bit suffix.  After resource evaluation, value
records carry the $b$-bit field value $G_y(z)$.  Both payloads fit within the
polynomial record-width bound.

\begin{lemma}[Record routing]\label{lem:routing}
Let a fixed table contain $M$ distinct keys.  Suppose every incidence key
belongs to this table and occurs in at most one of the $t(q-1)$ incidence
records.  With keys and payloads of width polynomial in
$n$, $\ell\log q$, $\log(t+2)$, and $\log(M+2)$, suffixes can be scattered
to their keyed slots and resource values gathered back to their request
records using circuits of total size
\[
  \widetilde O_\ell(tq+M).
\]
\end{lemma}

\emph{Implementation idea.}
The four passes in the table use a constant number of sorting networks on
$O(tq+M)$ records of polynomial width.  The exact fixed-wire treatment
of unmatched records and array boundaries is given in
\cref{app:routing-proof}.

For the resource bank above, the table consists of the $M=q^\ell$ field points.

\section{The composition bound}\label{sec:master}

Disjoint recovery lets each resource circuit serve at most one suffix.
There are $q^\ell b$ Boolean resources, each on $d$ variables, so ordinary
one-copy synthesis bounds their total size by $q^\ell bL(d)$.  The remaining
cost comes from selecting recovery sets, routing suffixes and values, and
XOR decoding.

\begin{proposition}[Scheduler-sensitive composition]
\label{prop:scheduler-sensitive}
Fix $\ell\ge2$ and the local-recovery gadget for a split $n=p+d$.
Suppose a Boolean circuit of size at most $S(t,q)$ maps every ordered list
of $t$ information points to enumerated pairwise disjoint affine-line
recovery sets.  Then every $f:\bits^n\to\bits$ satisfies
\begin{equation}\label{eq:scheduler-sensitive}
\C(f^{\times t})
\le \underbrace{q^\ell b\,L(d)}_{\text{resource circuits}}
   +\underbrace{S(t,q)}_{\text{schedule selection}}
   +\underbrace{\widetilde O_\ell\!\left(tq+q^\ell\right)}_{
     \text{routing and decoding}}.
\end{equation}
\end{proposition}

\begin{proof}
Map each prefix to its information point and bit position, and apply the
scheduler.  This costs $\widetilde O_\ell(t)+S(t,q)$ gates.  By
\cref{lem:routing}, scatter the assigned suffixes to the $q^\ell$ resource
positions, filling unused positions with a fixed dummy suffix.  Disjointness
ensures that each position receives at most one requested suffix.

Evaluate the $b$ Boolean functions $H_{y,r}$ at each position $y$, for a
total of $q^\ell bL(d)$ gates, and gather the resulting field values.
Routing costs $\widetilde O_\ell(tq+q^\ell)$.  For each request, XOR its
$q-1$ returned values and select the requested bit.  This is correct by
\cref{lem:line} and costs $O(tqb)$, which the displayed overhead absorbs.
\end{proof}

\section{Direct proof of exponential-range mass production}\label{sec:direct}

The parameter choice balances storage against recovery capacity.  Keeping
$d\approx\delta n$ suffix bits leaves about $2^{(1-\delta)n}$ encoded
prefix bits.  Since a recovery set occupies about $q$ points while a target
has about $q^{\ell-1}$ directions, the scheduler needs $t$ to be smaller
than a constant multiple of $q^{\ell-2}$.  The inequality below expresses
exactly this requirement at the level of exponential rates.

\begin{proof}[Direct proof of \cref{thm:main}]
Fix $\gamma<1$, choose any fixed $0<\delta<1-\gamma$, and put
$d=\lceil\delta n\rceil$ and $p=n-d$.  Choose a fixed $\ell\ge2$ large
enough that
\begin{equation}\label{eq:direct-margin}
  \gamma<(1-2/\ell)(1-\delta).
\end{equation}
Use the product-code gadget with its minimal packing field.  Its parameters
satisfy
\[
  \log_2q=(1-\delta)n/\ell+o(n),\qquad
  \log_2D_q=(1-1/\ell)(1-\delta)n+o(n).
\]
Thus \eqref{eq:direct-margin} implies $512tq\le D_q$ for all sufficiently
large $n$ and every $t\le2^{\gamma n}$.  Apply
\cref{thm:linear-scheduler,prop:scheduler-sensitive}.  This gives
\[
  \C(f^{\times t})\le q^\ell bL(d)
                 +\widetilde O_\ell(tq+q^\ell).
\]
The first term is $O_\ell(2^p2^d/d)=O_\gamma(2^n/n)$.  The two overhead
exponents are at most $\gamma+(1-\delta)/\ell<1$ and $1-\delta<1$,
respectively, with fixed positive margins.  All suppressed polynomial factors
are therefore negligible compared with $2^n/n$.  Enlarge the constant to
cover the finitely many initial lengths by ordinary circuit replication.
\end{proof}

\section{High-rate codes and the leading coefficient}\label{sec:high-rate}

The preceding proof has two constant losses.  First, the product code's rate
approaches $\ell^{-\ell}$: it stores substantially more symbols than the
information it carries.  Second, rounding up to the next allowed field size
can leave a constant fraction of the information positions unused.
We remove the first loss with a code of rate tending to one and the second
with several smaller codewords, so that only the last one needs padding.
The scheduler and decoder continue to use the same line-recovery interface.

Lifted Reed--Solomon codes already provide
rate tending to one with line-parity recovery
\cite{GuoKoppartySudan2013,HolzbaurEtAl2020}.  We give an elementary sufficient
monomial construction below, followed by the packing and circuit-cost
accounting.

\begin{lemma}[High-rate line-parity code]\label{lem:high-rate-code}
Fix integers $\ell\ge2$ and $h\ge\lceil\log_2\ell\rceil$.  Put $q=2^b$ and
$N=q^\ell$.  There is a
linear code on $\F_q^\ell$ of dimension $K$ such that
\begin{equation}\label{eq:high-rate}
  \frac KN\ge1-(1-2^{-\ell h})^{\lfloor b/h\rfloor}
       =1-o_{\ell,h}(1)\qquad(b\to\infty),
\end{equation}
and each word $P$ satisfies
\[
  P(u)=\sum_{\lambda\in\F_q^*}P(u+\lambda v)
\]
for every target $u$ and every nonzero direction $v$.  The code can be made
systematic on a fixed set of $K$ evaluation points.
\end{lemma}

\begin{proof}
The low-degree condition in \cref{lem:line} was sufficient for line parity,
but stronger than necessary.  For the sum over $\F_q$ to vanish, it is enough
that every positive degree appearing on a line avoid multiples of $q-1$.
We enforce this weaker condition using a common block of zero binary digits.

Consider reduced monomials $X_1^{e_1}\cdots X_\ell^{e_\ell}$ with
$0\le e_i<2^b$.  Partition the first $h\lfloor b/h\rfloor$ binary digit
positions into disjoint blocks of length $h$.  Retain a monomial if some block
is zero in \emph{all} of its exponents.  The fraction retained is exactly
$1-(1-2^{-\ell h})^{\lfloor b/h\rfloor}$.  These monomials are linearly
independent as evaluation functions, by repeated univariate interpolation.

To check line parity, expand a retained monomial at $u+\lambda v$.  In
characteristic two, every possibly nonzero term in this expansion has
$\lambda$-degree $J=\sum_i j_i$, where the one-bits of $j_i$ are a subset of
those of $e_i$.  This follows directly by writing
$(u_i+\lambda v_i)^{e_i}$ as a product over the one-bits of $e_i$ and using
$(a+c)^{2^r}=a^{2^r}+c^{2^r}$.

Suppose the common zero block starts at position $s$.  Every $j_i$ has the
same zero block, so its residue modulo $2^{s+h}$ is at most $2^s-1$.
Consequently,
\[
  J\bmod 2^{s+h}=\sum_i(j_i\bmod2^{s+h})
       \le\ell(2^s-1)\le2^{s+h}-\ell.
\]
There is no wraparound in this equality.  If $J$ were a positive multiple
$a(2^b-1)$, we would have $1\le a\le\ell-1$: the zero block makes every
$j_i<2^b-1$.  Since $s+h\le b$ and $a<2^{s+h}$, its residue would instead
be $2^{s+h}-a\ge2^{s+h}-\ell+1$, a contradiction.  Thus every positive
$\lambda$-degree is not divisible by $q-1$.  The monomial-sum calculation in
\cref{lem:line} gives zero for the sum over $\F_q$, including the constant
term.  Linearity proves line parity for the span of the retained monomials.

Finally, choose $K$ independent coordinate columns of a generator matrix.
Evaluation on those coordinates is an isomorphism to $\F_q^K$, giving a
systematic information set.  The choice and any change of basis are offline
parts of the nonuniform code specification.
\end{proof}

To index the information set, store its ordered list of $K$ points as a
hardwired table, and sort its records together with all $t$ live symbol
indices, placing a table record
before its queries at an equal index.  Propagate the table value along each
equal-index run, then sort queries back to request order.  This costs
$\widetilde O_\ell(K+t)$ gates and handles repeated indices.  It is the same
batched lookup primitive used for occupancy broadcast, with a point as payload.

\begin{proof}[Proof of \cref{thm:coefficient}]
Fix $0<\delta<1-\gamma$, put $d=\lceil\delta n\rceil$, $p=n-d$, and choose
a fixed $\ell$ satisfying \eqref{eq:direct-margin}.  Fix
$h=\lceil\log_2\ell\rceil$ in \cref{lem:high-rate-code}.  Instead of the
minimal packing field, take
\[
  b=\left\lfloor\frac{p-3\log_2(n+2)}{\ell}\right\rfloor,
  \qquad q=2^b,
  \qquad B=\left\lceil\frac{2^p}{Kb}\right\rceil,
\]
for sufficiently large $n$, where $K$ is the dimension in
\cref{lem:high-rate-code}.  Use $B$ separate copies of this code.  Each has
$Kb$ information bits, so they hold the $2^p$ prefix values with only the
last code's information positions padded.  Because $K/N\to1$ and
$Nb=o(2^p)$, their total number of Boolean resources is
\begin{equation}\label{eq:sharp-packing}
  2^p\le BNb
       \le\frac NK\,2^p+Nb=(1+o(1))2^p.
\end{equation}
Here $B=\poly_\ell(n)$ and
$\log_2q=(1-\delta)n/\ell+O_\ell(\log n)$.

Fixed-divisor arithmetic maps each prefix to its code number, information
index, and bit position.  The batched table lookup above supplies all target
points.  Schedule these $t$ targets together, even when they refer to different
code copies.  By \cref{thm:linear-scheduler}, the resulting recovery sets are
globally disjoint.  Attaching the requested code number to each incidence
therefore gives unique keys $(\text{code},y)$.  Route the $t(q-1)$ incidences
against $BN$ slots by the four-pass construction of \cref{lem:routing}.

For each of the $BNb$ Boolean resource functions on $d$ variables, use the
sharp one-copy synthesis bound
$L(d)\le(1+o(1))2^d/d$ in the stated basis
\cite{FrandsenMiltersen2005}.  The complete bound is
\[
  \C(f^{\times t})
  \le BNbL(d)+\widetilde O_\ell(tq+BN+K)
  \le\left(\frac1\delta+o(1)\right)\frac{2^n}{n}.
\]
Indeed, \eqref{eq:sharp-packing} controls the resource term; the other terms
have exponents at most $\gamma+(1-\delta)/\ell<1$ and $1-\delta<1$.
All estimates are uniform over $f$ and the allowed $t$.

For every fixed $\delta<1-\gamma$ we have proved a limiting upper coefficient
$1/\delta$.  Taking the infimum of these bounds as $\delta$ increases to
$1-\gamma$ proves the theorem.  This last step takes a limit of fixed-parameter
bounds; it makes no assertion uniform in growing $\ell$ or in $\gamma\to1$.
\end{proof}

\section{Linear depth at the same size order}\label{sec:linear-depth}

To prove \cref{thm:linear-depth}, we return to the constant-rate product
code of \cref{sec:code} and adapt its packing to make the prefix map a
wiring operation.  The scheduler also needs a depth refinement: halving
uses $O(\log t)$ dependent phases.  Testing larger menus lets each phase
leave only an exponentially small fraction of requests unfinished.
A constant number of phases then suffices, with total gate count still
negligible beside $2^n/n$.

We apply \cref{lem:collision-tail} with a smaller allowed remainder.
To implement a phase at low depth we use
Holmgren and Rothblum's sorting theorem
\cite[Lemma~6]{HolmgrenRothblum2024}: $R$ records of $w$ bits, with
$\kappa$-bit keys, can be sorted by a bounded-fan-in Boolean circuit of size
$O(Rw\log(R+2))$ and depth $O(\kappa+\log(R+2))$.
Their proof uses the sorting networks of Ajtai, Koml\'os, and Szemer\'edi,
with a bit-position argument that bounds the depth of the complete Boolean
network \cite[Section~4]{HolmgrenRothblum2024}.

\subsection{Menus that leave few requests unfinished}

\begin{lemma}[Menus with a prescribed remainder]\label{lem:small-remainder-menus}
Put $N=q^\ell$ and $E_g=g(1+3\log_2N)$.  Suppose $0<\alpha\le1/2$ and
\[
  \frac{4e\,gq}{D_q}\le\alpha^2,
\]
where $e$ is the base of the natural logarithm.
For each $1\le k\le g$, define
\[
  v_k=\lfloor\alpha k\rfloor,\qquad
  h_k=\left\lceil\frac{v_k+1}{2}\right\rceil,\qquad
  r_k=\left\lceil
       \frac{E_g+1}{h_k\log_2(1/\alpha)}
      \right\rceil.
\]
There is a fixed menu of $r_k$ direction $k$-tuples such that, for every
occupied set $U$ that is a union of at most $g$ recovery sets and every
ordered list of $k$ targets, a menu entry has at least $k-v_k$ clean
recovery sets.  Moreover,
\begin{equation}\label{eq:small-remainder-menu-work}
  r_k k\le
  g+\frac{2(E_g+1)}{\alpha\log_2(1/\alpha)}.
\end{equation}
\end{lemma}

\begin{proof}
We have $h_k>\alpha k/2$.  By \cref{lem:collision-tail} and the inequality
$\binom{k}{h}\le(ek/h)^h$, one random tuple fails with probability at most
\[
  \left(\frac{ek}{h_k}\frac{2gq}{D_q}\right)^{h_k}
  \le
  \left(\frac{4e\,gq}{\alpha D_q}\right)^{h_k}
  \le\alpha^{h_k}.
\]
As in \cref{lem:universal-menus}, there are at most
$(1+ND_q)^gN^k\le2^{E_g}$ states to serve.  A menu of $r_k$ independent
tuples fails on any particular state with probability at most
$2^{-r_kh_k\log_2(1/\alpha)}\le2^{-E_g-1}$.
The union bound is less than one, so a menu works for all states.
Finally, the ceiling in $r_k$ and $k/h_k<2/\alpha$ give
\eqref{eq:small-remainder-menu-work}.
\end{proof}

In a successful phase, accept exactly the first $k-v_k$ clean requests,
reserve their recovery sets, and compact the remaining $v_k$ requests.
The successive values of $k$ are fixed in advance, since each is replaced
by $\lfloor\alpha k\rfloor$.  The number of phases is at most
\[
  1+\left\lfloor\frac{\log_2g}{\log_2(1/\alpha)}\right\rfloor.
\]
This includes the final phase, when $\lfloor\alpha k\rfloor=0$.

\subsection{Depth of one phase and of the resources}

For a phase with $k$ pending requests, generate the $r_k k(q-1)$
candidate-point records and append the occupied points once.  Thus the
record count is
\[
  R=O\bigl(q(g+r_k k)\bigr).
\]
The three verification passes from \cref{thm:linear-scheduler} still
apply.  Candidate directions and line scalars are hardwired, so each
candidate point is its target plus a fixed field vector.

The occupancy propagation in the first pass need not be a serial scan.
Writing its update as $e_i=a_i\vee(b_i\wedge e_{i-1})$, encode it by
the pair $(a_i,b_i)$.  Composition is the associative operation
\[
  (a,b)\circ(c,d)=\bigl(a\vee(b\wedge c),\,b\wedge d\bigr).
\]
A doubling prefix circuit therefore computes every occupancy bit in
$O(\log(R+2))$ depth and $O(R\log(R+2))$ gates.  This is the usual
parallel-prefix method; a related propagation primitive is given in
\cite[Lemma~26]{HolmgrenRothblum2024}, with credit there to Ladner and
Fischer.

OR the bad-point flags along each request's line with a balanced tree.
To test whether a candidate has at least $k-v_k$ clean requests, sort
its $k$ clean bits and inspect the appropriate fixed position.
Prefix OR selects the first successful candidate.  One-hot multiplexers
select its records, and sorting by the clean flag and request position
accepts exactly the required number of clean requests.  All these
operations have depth $O(w+\log(R+2))$ and size
$R\,\operatorname{poly}(w+\log(R+2))$, when all keys and records have
width $O(w)$.  The record sorting steps use
\cite[Lemma~6]{HolmgrenRothblum2024}.  In the parameter choice below,
$w=O_\gamma(n)$ and $\log(R+2)=O_\gamma(n)$, so a phase has depth
$O_\gamma(n)$.

The resource circuits run in parallel, so their contribution to depth is
the maximum depth of any one resource.  The following synthesis bound
keeps this part linear as well.

\begin{lemma}[One-copy synthesis with a depth bound]\label{lem:size-depth-synthesis}
Every Boolean function on $d\ge1$ bits has a circuit of size $O(2^d/d)$
and depth $d+O(\log(d+2))$.
\end{lemma}

\begin{proof}
For $d\ge4$, let $k=\lfloor\log_2d\rfloor-1$ and split the input into
$d-k$ prefix bits and $k$ suffix bits.  Compute all minterms of the prefix
with $O(2^{d-k})$ gates and depth $O(\log(d+2))$: recursively form the
minterms on the two halves and AND every pair of half-minterms.

On the suffix, compute all $2^{2^k}$ Boolean functions.  Computing all
suffix minterms once and then the OR for each truth table uses
$O(2^k2^{2^k})$ gates and depth $k+O(\log(k+2))$.
For each prefix assignment, select by fixed wiring the suffix function
specified by the desired restriction.  AND it with that prefix's
minterm, and OR all $2^{d-k}$ terms using a balanced tree.
The size is
\[
  O(2^{d-k}+2^k2^{2^k})=O(2^d/d),
\]
because $d/4<2^k\le d/2$.  The final OR contributes $d-k$ layers, giving
the stated depth.  Supplying negated input bits adds at most one layer.
The finitely many smaller values of $d$ are absorbed in the constants.
\end{proof}

\subsection{The simultaneous bound}

\begin{proof}[Proof of \cref{thm:linear-depth}]
Fix $0\le\gamma<1$, choose $0<\delta<1-\gamma$, and take a fixed
$\ell$ sufficiently large that
\[
  \gamma<(1-2/\ell)(1-\delta).
\]
Put $d=\lceil\delta n\rceil$, $p=n-d$, and
$\lambda=(1-\delta)/\ell$.  We use the product-code construction, with
a minor packing choice that makes the prefix map a wiring operation.

Let $c=\lceil\log_2\ell\rceil+1$.  For a sufficiently large field
$q=2^b$, take $A$ to be the first $2^{b-c}$ elements in the fixed binary
representation of $\F_q$.  Use the product grid $A^\ell$, of size
$K=2^{\ell(b-c)}$.  The interpolating polynomial has degree at most
$\ell(2^{b-c}-1)\le q-2$, so the same line-parity recovery identity
applies.  Pack only $b'=2^{\lfloor\log_2 b\rfloor}$ information bits
into each field symbol, padding its remaining $b-b'$ bits with zeros.
Choose the least sufficiently large $b$ such that $Kb'\ge2^p$.

Since $Kb'$ is a power of two, pad a prefix with leading zero bits
and split it into $\ell$ coordinate indices of $b-c$ bits and one
bit-position index of $\log_2b'$ bits.  Prepending $c$ zeros to each
coordinate index gives its field representation.  Thus no division or
table lookup is needed for this prefix map.
Increasing $b$ by one multiplies $Kb'$ by at most $2^{\ell+1}$, while
$b<2b'$.  Minimality therefore gives
\[
  q^\ell b=O_\ell(2^p),\qquad
  \log_2q=\lambda n+O_\ell(\log(n+2)).
\]
The constant loss in this packing is harmless for the order-of-growth
bound.

Set
\[
  \beta=(\ell-2)\lambda-\gamma>0,\qquad
  \rho=\beta/4,\qquad
  \alpha=2^{-\lfloor\rho n\rfloor}.
\]
For all sufficiently large $n$ and uniformly for $g=t\le2^{\gamma n}$,
\[
  \frac{gq}{D_q}\le2^{-\beta n+O_\ell(\log(n+2))},
\]
so $4e\,gq/D_q\le\alpha^2$ and $0<\alpha\le1/2$.
Apply \cref{lem:small-remainder-menus}.
Since $E_g=O_\gamma(gn)$ and $\log_2(1/\alpha)=\Theta_\gamma(n)$,
\eqref{eq:small-remainder-menu-work} gives $r_k k=O_\gamma(g/\alpha)$
uniformly over every phase.  The number of phases is $O_\gamma(1)$,
and each phase uses
\[
  R\le C_\gamma gq/\alpha
  \le2^{(\gamma+\lambda+\rho)n+O_\gamma(\log(n+2))}
\]
records per phase, for a suitable constant $C_\gamma$.  Moreover,
\[
  \gamma+\lambda+\rho
  =\tfrac34\gamma+\tfrac{\ell+2}{4}\lambda
  <(\ell-1)\lambda<1.
\]
The scheduler consequently has depth $O_\gamma(n)$ and size
$o_\gamma(2^n/n)$, even with the polynomial factors from its
Boolean implementation.  Its menus remain fixed nonuniform data.

Scatter and gather use the four passes of \cref{lem:routing}, now
implemented with Holmgren--Rothblum sorting.  There are
$O(gq+q^\ell)$ records of $O_\gamma(n)$ bits.  Each pass has depth
$O_\gamma(n)$, and their total size is $o_\gamma(2^n/n)$ because
$\gamma+\lambda<1$ and $\ell\lambda=1-\delta<1$.
Balanced XOR trees recover each field symbol in $O(\log q)$ depth;
selecting its requested bit adds $O(\log(b+2))$ depth.

Finally, apply \cref{lem:size-depth-synthesis} independently to the
$q^\ell b$ Boolean resource functions.  Their total size is
\[
  O\left(q^\ell b\,\frac{2^d}{d}\right)
  =O_\gamma(2^n/n),
\]
and they run in parallel at depth $O_\gamma(n)$.
The prefix map, the constant number of scheduling phases, routing,
resource evaluation, and decoding therefore have total depth
$O_\gamma(n)$ and total size $O_\gamma(2^n/n)$.
Enlarge both constants to handle the finitely many initial input lengths.
\end{proof}

The power-of-two packing and \cref{lem:size-depth-synthesis} allow
constant losses, so this proof controls the size order without retaining
the leading coefficient of \cref{thm:coefficient}.

\subsection{Randomized construction from the truth table}

Amplifying the menu union bound turns the preceding existence proof into
a randomized construction.

\begin{corollary}[Randomized construction]\label{cor:randomized-construction}
For every fixed $0\le\gamma<1$, there is a randomized algorithm that,
given the truth table of $f:\bits^n\to\bits$ and an integer
$1\le t\le2^{\gamma n}$, runs in time polynomial in $2^n$ and outputs
a circuit of size $O_\gamma(2^n/n)$ and depth $O_\gamma(n)$.
With probability at least $1-2^{-n}$ over the construction, that circuit
computes $f^{\times t}$ correctly on every input batch.
\end{corollary}

\begin{proof}
Choose a rational rate $\gamma<\widehat\gamma<1$ and use the proof of
\cref{thm:linear-depth} at rate $\widehat\gamma$, also choosing rational
$\delta$.  Its fixed numerical parameters are then rational, and its
bounds remain $O_\gamma(2^n/n)$ and $O_\gamma(n)$.
Let $P=O_\gamma(1)$ be its number of phases.  The successive pending
request counts $k$ are fixed before any input is supplied.  In each phase,
sample a menu of independent uniform direction tuples, replacing its
length by
\[
 r_k'=\left\lceil
 \frac{E_g+n+\lceil\log_2P\rceil+2}
      {h_k\log_2(1/\alpha)}
 \right\rceil.
\]
The proof of \cref{lem:small-remainder-menus} bounds the probability that
this menu fails on any state by
\[
 2^{E_g-r_k'h_k\log_2(1/\alpha)}
 \le 2^{-n-\lceil\log_2P\rceil-2}.
\]
A union bound over the phases is at most $2^{-n-2}$.  On the
complementary event all menus work for every state, hence for every
input batch and every choice of the resource truth tables.  Since
$E_g=\Theta_\gamma(gn)$, the enlarged menus still satisfy
$r_k'k=O_\gamma(g/\alpha)$, preserving both circuit bounds.

All remaining construction steps take polynomial time in $2^n$.
The field has $b=O_\gamma(n)$ bits per element; exhaustive search and
trial division find an irreducible polynomial of degree $b$ in
$2^{O_\gamma(n)}$ time.  Product-grid interpolation constructs
the resource truth tables in polynomially many field operations in their
total table length, which is $2^{O_\gamma(n)}$.  The synthesis in
\cref{lem:size-depth-synthesis} explicitly lists the required gates.
The sorting circuits are log-space uniform
\cite[Lemmas~8--9]{HolmgrenRothblum2024}, and the other scheduling and
routing gates are explicitly specified above.

For a worst-case polynomial running time, implement each uniform
direction choice by rejection sampling an index in $[D_q]_0$.
If $L$ directions are needed in total, cap each choice at
$n+\lceil\log_2(L+1)\rceil+2$ trials.  Each trial succeeds with probability
at least $1/2$, so the probability of any exhausted cap is at most
$2^{-n-2}$.  On that event output the all-zero circuit.  Otherwise the
chosen directions remain independent and uniform.  Combining the two
failure bounds proves the stated probability.  The finitely many initial
input lengths use elementary truth-table synthesis, with constants enlarged.
\end{proof}

The output circuit uses no runtime randomness.  The algorithm does not
certify that its sampled menus work, and its time bound is polynomial
in the full truth table, which may be exponentially larger than a
succinct description of $f$.


\section{Optimal approximate quantum state mass production}
\label{sec:quantum}

The upper bound in \cref{thm:quantum} has three steps: evaluate Boolean
functions coherently, use their outputs to select rotation angles, and
assemble the target state from conditional rotations.  The lower bound
compares a packing of tensor-power states with the number of finite-gate
circuits.  Throughout, we charge every elementary gate and measure error
on the entire batch.

\subsection{Model and coherent Boolean evaluation}

Fix the finite gate set
$\mathcal G=\{H,T,T^\dagger,\mathrm{CNOT}\}$, with arbitrary connectivity.
Here $H$ is the Hadamard gate, $T=\operatorname{diag}(1,e^{i\pi/4})$,
and CNOT is the controlled-NOT gate.  Every gate has cost one; quantum
depth counts layers of gates acting on pairwise disjoint qubits.
Circuits have no measurements or postselection.

For state preparation, all qubits start in zero, the $nt$ output qubits
are fixed, and every scratch qubit returns exactly to zero.  The number
of scratch qubits is unrestricted, and the circuit may depend on the
target state.  The coherent subroutines below instead act on arbitrary
input states, with the same requirement on scratch qubits.

For normalized pure states, write
\[
 d(\psi,\phi)=\sqrt{1-|\langle\psi\mid\phi\rangle|^2}
\]
for their trace distance; this identifies states differing only by a
global phase.  Let $Q_\varepsilon(\psi,t)$ be the minimum gate count
whose output is within trace distance $\varepsilon$ of
$|\psi\rangle^{\otimes t}$, and set
\[
 M(n,t)=\max_{\psi\in\mathbb C^{2^n},\ \|\psi\|=1}
                   Q_{1/10}(\psi,t).
\]
The construction below bounds these integer-valued minima uniformly,
so the maximum is well-defined.  The state is specified when the circuit
is designed; no copy of it is supplied as an input.

We first implement Boolean mass production reversibly.  The simulation
also gives a depth bound for the resulting oracle.

\begin{lemma}[Coherent Boolean mass production]\label{lem:coherent}
If a Boolean circuit of size $s$ and depth $D$ computes $f^{\times t}$,
there is a quantum circuit of size $O(s+t+1)$, depth
$O((D+1)\log(s+t+2))$, and $O(s+t+1)$ scratch qubits implementing
\[
 |x_1,\ldots,x_t\rangle|z_1,\ldots,z_t\rangle|0^a\rangle
 \longmapsto
 |x_1,\ldots,x_t\rangle
 |z_1\xor f(x_1),\ldots,z_t\xor f(x_t)\rangle|0^a\rangle.
\]
In particular, for fixed $0\le\eta<1$, $k\ge1$, and
$t\le2^{\eta k}$, this map has size $O_\eta(2^k/k+t)$ for every
$k$-bit Boolean function $f$.
\end{lemma}

\begin{proof}
Delete gates outside the output cones.  Assign a separate control qubit
to each use of a Boolean wire, including its uses as a designated output.
A wire with $r$ uses supplies these controls by a balanced CNOT tree,
using $O(r)$ gates and scratch qubits
at depth $O(\log(r+1))$.  The total number of uses is at most $2s+t$,
so this costs $O(s+t+1)$ gates and scratch qubits in all.

Compute the Boolean gates in depth order, each into a fresh zero qubit.
NOT, AND, and OR each require a constant number of NOT, CNOT, and
Toffoli gates; repeated sources and constants are handled by their
one-input or constant simplifications.  These gates have exact
constant-size implementations over $\mathcal G$.  The separate control
qubits make the simulations within a Boolean layer disjoint.  After
each layer, prepare its output fan-out trees in parallel.  The initial
input fan-out and the $D$ gate layers therefore have total depth
$O((D+1)\log(s+t+2))$.

Copy the designated output controls to the distinct $z$ qubits in
parallel, then reverse all gate computations and fan-out trees.
Prepared constant sources and all scratch qubits are restored.

The displayed equality holds on every computational basis state, hence
by linearity on every superposition and on inputs entangled with other
registers.  Apply \cref{thm:main} for the final assertion.
\end{proof}

\begin{corollary}[Coherent evaluation at quadratic depth]
\label{cor:coherent-depth}
For fixed $0\le\gamma<1$, every $n$-bit Boolean function and every
$1\le t\le2^{\gamma n}$ admit the exact coherent evaluation map in
\cref{lem:coherent} with $O_\gamma(2^n/n)$ elementary gates,
$O_\gamma(n^2)$ depth, and $O_\gamma(2^n/n)$ scratch qubits, all restored
to zero.
\end{corollary}

\begin{proof}
Apply \cref{lem:coherent} to \cref{thm:linear-depth}, using
$t=O_\gamma(2^n/n)$ and $\log(s+t+2)=O_\gamma(n)$.
\end{proof}

For state preparation we use the gate bound in \cref{lem:coherent}.
The remaining rotation and synthesis steps below are analyzed for size only.

\subsection{Conditional rotations with clean scratch space}

For $a\in\{y,z\}$, let $R_a(\theta)=\exp(-i\theta\sigma_a/2)$,
where $\sigma_a$ is the corresponding Pauli matrix.  A \emph{multiplexed
rotation} on $k$ address qubits and one target qubit is
\[
 U_\theta=\sum_{x\in\bits^k}|x\rangle\langle x|
                           \otimes R_a(\theta(x)).
\]
Thus an address selects a rotation angle while remaining unchanged.
We use $\|\cdot\|$ for the operator norm on operators.

\begin{lemma}[Batched conditional rotations]\label{lem:quantum-rotation}
Fix $0\le\eta<1$.  Suppose $k\ge1$, $t\le2^{\eta k}$, and
$0<\delta\le1/4$.  Put $b=\lceil\log_2(8\pi/\delta)\rceil$.
For every multiplexed rotation $U_\theta$, a circuit over $\mathcal G$
with clean scratch space approximates $U_\theta^{\otimes t}$ to
operator-norm error at most $t\delta$, using
\[
 O_\eta\!\left(b\frac{2^k}{k}+t\,\poly(b)\right)
\]
gates.  The polynomial is fixed, independent of $k,t,\theta$.
The circuit preserves every address qubit exactly.
\end{lemma}

\begin{proof}
Represent each angle modulo $4\pi$, and round it down to
$4\pi\sum_{j=1}^b2^{-j}f_j(x)$, where each $f_j$ is a Boolean
function of the address.  Since
$\|R_a(u)-R_a(v)\|\le |u-v|/2$, this changes each selected rotation
by at most $2\pi2^{-b}\le\delta/4$.

For each $j$, use \cref{lem:coherent} to compute
$f_j(x_1),\ldots,f_j(x_t)$ into zero qubits.  On block $i$, rotate
the target by $4\pi2^{-j}$ controlled by the computed bit, and then
uncompute the bits.  This charges
$O_\eta(b2^k/k+bt)$ gates for Boolean evaluation and its reversal.
The address registers and all computed control bits are unchanged by
the controlled rotations, so reversing the lookup erases its scratch
space exactly.

It remains to synthesize a controlled rotation in the finite gate set
without losing this property.  Write $X$ for the NOT matrix.  Given a
one-qubit word $A$ approximating $R_a(\theta/2)$, use the two-CNOT
circuit with block form
\[
 |0\rangle\langle0|\otimes I+
 |1\rangle\langle1|\otimes A X A^\dagger X.
\]
The identity $X R_a(u)X=R_a(-u)$ shows that the ideal choice of $A$
implements the desired controlled rotation.  For every implemented
word $A$, the control-zero block is exactly $I$ and the control is
preserved exactly.  A global phase in $A$ cancels against $A^\dagger$.
The Solovay--Kitaev theorem \cite{DawsonNielsen2006} supplies $A$ of
length polynomial in $\log(b/\delta)$ so that this controlled gate
has error at most $\delta/(4b)$.

After lookup and uncomputation the resulting operation on the data is
the tensor product of the corresponding approximate one-block
multiplexors: each computed bit depends only on its own block's address.
The error in one block is at most $\delta/4+b\delta/(4b)\le\delta$.
The tensor-product telescoping inequality gives error at most $t\delta$.
There are $bt$ synthesized controlled gates.  Since
$\log(b/\delta)=O(b)$, their total cost is $t\poly(b)$, as claimed.
\end{proof}

\begin{lemma}[State preparation by conditional rotations]
\label{lem:state-rotations}
Up to global phase, every $n$-qubit pure state can be prepared using
two multiplexed rotations, one about each of the $y$ and $z$ axes,
at each address length $k=0,\ldots,n-1$.
\end{lemma}

\begin{proof}
This is the standard recursive state-synthesis decomposition
\cite{ShendeBullockMarkov2006,Kretschmer2023}.  To see it directly,
write the target as
\[
 |\psi\rangle=\sum_{x\in\bits^{n-1}}|x\rangle
                   (a_x|0\rangle+b_x|1\rangle).
\]
For each $x$, choose angles $\theta_x,\phi_x$ and a complex scalar
$c_x$ such that
\[
 a_x|0\rangle+b_x|1\rangle
   =c_x R_z(\phi_x)R_y(\theta_x)|0\rangle,
 \qquad |c_x|^2=|a_x|^2+|b_x|^2.
\]
Two rotations parametrize any normalized one-qubit state up to phase;
when both amplitudes vanish take $c_x=0$ and arbitrary angles.
The state $\sum_x c_x|x\rangle$ is normalized.  Prepare it recursively,
adjoin a zero target, and apply the two rotations controlled by $x$.
The induction ends with a scalar of modulus one, which contributes
only an irrelevant global phase.
\end{proof}

\subsection{Upper bound and precision allocation}

To control the error of the whole batch, we give each copy an $O(1/t)$
error budget.  We allocate most of that budget to the largest address
tables, which then need only $O(\log(t+1))$ angle bits.  Smaller tables
use finer precision at a geometrically smaller synthesis cost.  Related
allocations appear in \cite[Appendix~B]{GossetKothariWu2026}.

\begin{proposition}[Quantum upper bound]\label{prop:quantum-upper}
For every fixed $0\le\gamma<1$, every $n$-qubit pure state $\psi$,
and every $1\le t\le2^{\gamma n}$,
\[
 Q_{1/10}(\psi,t)
   =O_\gamma\!\left(\frac{2^n}{n}\log_2(t+1)\right).
\]
\end{proposition}

\begin{proof}
Fix constants $\gamma<\alpha<1$ and $\gamma/\alpha<\eta<1$,
and put $m=\lceil\alpha n\rceil$ and $L=1+\log_2t$.
For sufficiently large $n$ we have $1\le m<n$.
Use \cref{lem:state-rotations} on each of the $t$ blocks.  At address
length $k$, approximate each of the two rotations with the one-block
error allowance
\[
 \delta_k=\frac{1}{40t}\,2^{k-n}.
\]
The number of angle bits in \cref{lem:quantum-rotation} is then
\[
 b_k=O(n-k+L)=O(n).
\]
Since $t\delta_k$ bounds the error of each batched rotation, telescoping
over the stages gives total operator-norm error at most
\[
 2t\sum_{k=0}^{n-1}\delta_k
   =\frac1{20}\sum_{k=0}^{n-1}2^{k-n}<\frac1{10}.
\]
Acting on the all-zero state gives the same bound on vector distance
after aligning global phases, and hence on trace distance.

To apply the Boolean theorem at short address lengths, replace $k$ by
$k'=\max\{k,m\}$.  When $k<m$, add zero address qubits and extend
each angle function to ignore them.  They are preserved and can be
reused.  The fixed rate $\eta$ works at every stage, since
\[
 \log_2t\le\gamma n<\eta\alpha n\le\eta k'.
\]
Thus no stage invokes a mass-production bound outside its allowed range.

For $k\ge m$, \cref{lem:quantum-rotation} bounds the Boolean lookup
cost, summed over the two rotations per stage, by
\begin{align*}
 O_\gamma\!\left(\sum_{k=m}^{n-1}(n-k+L)\frac{2^k}{k}\right)
 &\le O_\gamma\!\left(\frac{2^n}{n}
                  \sum_{j\ge1}(j+L)2^{-j}\right)\\
 &=O_\gamma\!\left(\frac{2^n}{n}L\right).
\end{align*}
Here $k\ge\alpha n$ bounds the denominators, and the two geometric
sums are finite.  The $O(n)$ short stages each have $b_k=O(n)$ and
address length $m$, for a total lookup cost
\[
 O_\gamma(n^2 2^m/m)=O_\gamma(n2^{\alpha n})=o(2^n/n).
\]
Finally, all rotation-synthesis and output-copying costs sum to
$t\poly(n)=o(2^n/n)$ uniformly for $t\le2^{\gamma n}$.
The strict fixed margins $\alpha<1$ and $\gamma<1$ justify both
absorptions.  The scratch space is restored exactly at each stage.
Enlarging the constant covers the finitely many remaining input lengths.
\end{proof}

\subsection{Packing tensor powers and counting circuits}

Taking tensor powers makes nearby states easier to distinguish.  We
therefore pack one-copy states at distance of order $1/\sqrt t$; their
tensor powers remain separated by a constant.

\begin{lemma}[Packing tensor-power states]\label{lem:quantum-packing}
For $D\ge2$ and $t\ge1$, there is a set $\mathcal P$ of pure states
in $\mathbb C^D$ with $|\mathcal P|\ge(4t)^{D-1}$ such that distinct
$\psi,\phi\in\mathcal P$ satisfy
\[
 d(\psi^{\otimes t},\phi^{\otimes t})
     \ge\sqrt{1-e^{-1/4}}>2/10.
\]
\end{lemma}

\begin{proof}
For a fixed pure state $\psi$ and a Haar-random pure state $\phi$,
a trace-distance ball of radius $r\in[0,1]$ has probability
\[
 \Pr[d(\psi,\phi)\le r]=r^{2(D-1)}.
\]
Indeed, unitary invariance lets us take $\psi$ to be the first basis
vector.  Normalize $D$ independent standard complex Gaussian entries
to generate $\phi$.  Their squared moduli are independent exponentials,
and the first divided by their sum has the
$\operatorname{Beta}(1,D-1)$ distribution.  Its upper-tail probability
at $1-r^2$ is $(r^2)^{D-1}$.

Choose a maximal set of states separated by at least
$r=1/(2\sqrt t)$.  Compactness guarantees such a finite set, and
maximality makes its radius-$r$ balls cover the pure-state space.
Their volumes give $|\mathcal P|r^{2(D-1)}\ge1$.
For distinct members,
\begin{align*}
 d(\psi^{\otimes t},\phi^{\otimes t})^2
  &=1-|\langle\psi\mid\phi\rangle|^{2t}\\
  &\ge1-\left(1-\frac1{4t}\right)^t
   \ge1-e^{-1/4}.
\end{align*}
This proves both assertions.
\end{proof}

\begin{proposition}[Quantum lower bound]\label{prop:quantum-lower}
Uniformly over integers $1\le t\le2^n$ as $n\to\infty$,
\[
 M(n,t)=\Omega\!\left(nt+\frac{2^n}{n}\log_2(t+1)\right).
\]
\end{proposition}

\begin{proof}
Put $D=2^n$ and use \cref{lem:quantum-packing}.  The triangle
inequality implies that one circuit output can approximate at most
one of the $|\mathcal P|$ tensor-power targets to error $1/10$.

Let $g\ge1$ bound the number of gates.  A circuit uses at most $2g$
scratch qubits that participate in any gate; unused zero qubits can
be removed.  Canonically relabel the scratch qubits by their first use,
keeping the $nt$ designated output wires fixed.  A gate is specified
by its type in the fixed finite set $\mathcal G$ and at most two wire
indices in a set of size $nt+2g$.  Therefore the number of circuits
of size at most $g$ is at most
\[
 (g+1)\bigl(|\mathcal G|(nt+2g+2)^2\bigr)^g.
\]
This upper bound includes circuits that fail to restore their scratch
space, so imposing the clean-scratch condition cannot invalidate it.

If every target had a circuit of size at most $g$, taking logarithms
and comparing with $|\mathcal P|$ would give
\[
 (D-1)\log_2(4t)
    \le O\!\left(g\log_2(nt+g+2)+\log_2(g+1)\right).
\]
For $g\le2^n$ and $t\le2^n$, the logarithm on the right is $O(n)$,
which yields the logarithmic term.  If $g>2^n$, that term already follows
from $\log_2(t+1)\le n+1$.

For the remaining $nt$ term, take $|\psi\rangle=|1^n\rangle$.
Any untouched output wire stays in zero, making the output orthogonal
to $|1^{nt}\rangle$.  Error below one therefore requires every output
wire to participate in a gate, and a gate touches at most two wires.
Thus $M(n,t)\ge nt/2$.  Combining the two lower bounds proves the claim.
\end{proof}

\Cref{prop:quantum-upper,prop:quantum-lower} prove
\cref{thm:quantum}.  The packing argument gives a worst-case lower bound;
it does not identify an explicit family of hard states.


\section{Local hardness amplification and mass production}\label{sec:amplification}

Mass production can coexist with strong average-case circuit hardness.
We construct encoded functions for which even prediction with exponentially
small advantage requires almost $2^n/n$ gates, while exponentially many
exact evaluations cost $O(2^n/n)$ gates in total.  The construction combines
Hirahara's local hardness amplifier with our synthesis theorem.

Fix a universal machine and let $K(f)$ denote the plain, time-unbounded
Kolmogorov complexity of the truth table of $f$.  For $0\le\rho<1/2$, write
\[
 K_\rho(g)=\min_{h:\,\Pr_x[h(x)\ne g(x)]\le\rho}K(h),
 \qquad
 \C_\rho(g)=\min_{h:\,\Pr_x[h(x)\ne g(x)]\le\rho}\C(h),
\]
where $h$ has the same domain as $g$ and $x$ is uniform on that domain.

\begin{theorem}[Local amplification with an exponential batch]
\label{thm:amplification-batch}
For every fixed $0\le\tau<1$ and $\eta>0$, there exist an integer
$c\ge1$, a rational constant $\alpha>0$, and encodings $E_n$ taking
functions on $n$ bits to functions on $cn$ bits with the following
properties.  Write $g=E_n(f)$ and
$\varepsilon_n=2^{-\lfloor\alpha n\rfloor}$.  Uniformly over all $f$
and integers $1\le t\le2^{\tau n}$, for all sufficiently large $n$,
\begin{align}
 K_{1/2-\varepsilon_n}(g)&\ge K(f)-o(2^n),
       \label{eq:amplification-retention}\\
 \C(g^{\times t})&\le
    \left(\frac{1}{1-\tau}+\eta\right)\frac{2^n}{n}.
       \label{eq:amplification-batch-size}
\end{align}
There are also circuits for these batches of size
$O_{\tau,\eta}(2^n/n)$ and depth $O_{\tau,\eta}(n)$.
Given the source truth table, the full encoded table is computable
in time polynomial in $2^n$.
\end{theorem}

\begin{proof}
Hirahara's amplifier \cite[Lemma~8.1]{Hirahara2022}, based on
Impagliazzo--Wigderson \cite{ImpagliazzoWigderson1997}, supplies constants
$c\in\mathbb N$ and $a,b\ge1$ with the following properties.  For
rational parameters $\varepsilon,\delta\in(0,1/2)$, it encodes $f$ as a function
$g$ on $cn$ bits using a nonadaptive $f$-oracle circuit with
\[
 q=O(1/(\varepsilon\delta)),\quad
 s=O\!\left((n/(\varepsilon\delta))^a\right),\quad
 d=O(\log(n/(\varepsilon\delta))),
\]
queries, ordinary gates, and depth, respectively.  Taking its
reconstruction exponent to be $1/2$ gives
\begin{equation}\label{eq:amplifier-description}
 K_{1/2-\varepsilon}(g)
 \ge K(f)-O\!\left(2^{n/2}(1/(\varepsilon\delta))^b\right)
              -2^nH_2(\delta)-O(n),
\end{equation}
where $H_2$ is binary entropy.  The encoded table is computable from
the source table in time polynomial in $2^n$ and $1/(\varepsilon\delta)$.

Choose $\gamma\in(\tau,1)$ with
$1/(1-\gamma)<1/(1-\tau)+\eta/2$, and then a rational $\alpha>0$ such that
\begin{equation}\label{eq:amplification-margins}
 \tau+\alpha<\gamma,\qquad
 \tau+a\alpha<1,\qquad
 \frac12+b\alpha<1.
\end{equation}
Set $\varepsilon=\varepsilon_n$ and $\delta=1/n$.  For all sufficiently
large $n$, the total queries and ordinary gates for $t$ encoded evaluations
satisfy
\[
 tq=O\!\left(n2^{(\tau+\alpha)n}\right)\le2^{\gamma n},
 \qquad
 ts\le2^{(\tau+a\alpha)n}\poly(n)=o(2^n/n).
\]
Nonadaptivity lets us compute all addresses first and feed them to one
circuit for $f^{\times tq}$, then use its answers in the output circuits.
The total size is at most $ts+\C(f^{\times tq})$, even when addresses
repeat or share input bits.  \Cref{thm:coefficient} gives
\eqref{eq:amplification-batch-size}; \cref{thm:linear-depth} gives the
separate simultaneous size and depth bound, since the ordinary circuitry
before and after the oracle layer has total depth at most $2d=O(n)$.

The description loss in \eqref{eq:amplifier-description} is at most
\[
 2^{(1/2+b\alpha)n}\poly(n)+2^nH_2(1/n)+O(n)=o(2^n),
\]
proving \eqref{eq:amplification-retention}.  The margins in
\eqref{eq:amplification-margins} are fixed independently of $f$ and $t$,
so all estimates are uniform.  The encoding time is polynomial in $2^n$
because $1/(\varepsilon_n\delta)\le n2^{\alpha n}$.
The finitely many smaller $n$ may be assigned arbitrary encodings.
\end{proof}

\begin{corollary}[Approximation retains the source synthesis cost]
\label{cor:robust-synthesis}
For the encoding in \cref{thm:amplification-batch} and all sufficiently
large $n$, a uniformly random $f:\bits^n\to\bits$ satisfies,
with probability at least $1-2^{-n}$,
\[
 \C_{1/2-\varepsilon_n}(E_n(f))
       \ge (1-o(1))\frac{2^n}{n}.
\]
On the same event, all the exact batch upper bounds in that theorem
hold.  In particular, for $\tau=0$ and any fixed $\eta>0$, there
are functions $g$ with
\[
 (1-o(1))\frac{2^n}{n}
 \le\C_{1/2-\varepsilon_n}(g)
 \le\C(g)\le(1+\eta)\frac{2^n}{n}.
\]
\end{corollary}

\begin{proof}
With probability at least $1-2^{-n}$, counting programs gives
$K(f)\ge2^n-n$.  By \eqref{eq:amplification-retention}, every $h$
agreeing with $E_n(f)$ on at least a $1/2+\varepsilon_n$ fraction of
inputs then has $K(h)\ge(1-o(1))2^n$.  A smallest circuit for $h$ with
$S$ gates and $cn$ inputs has a description of length
\[
 K(h)\le(S+1)\log_2(S+cn+2)+O(S+\log(n+2))
\]
by \cref{lem:multi-output-count}, including the specification of $n,S$
and a fixed evaluation procedure.  For $S\le2^{n+1}/n$, this is at most
$nS+o(2^n)$, uniformly in $S$; larger $S$ already satisfies the desired
lower bound.  The upper bounds hold for every source.
\end{proof}

The size scale is the source length $L=2^n$, although the encoded table
has length $2^{cn}$.  The source is random, so this does not identify an
explicit family with exponential circuit complexity.  The advantage
exponent $\alpha$ may decrease as $\tau$ approaches one or $\eta$ approaches
zero, and the sharp coefficient and linear depth are separate guarantees.
Hardwiring the source into the multiselection circuits of
\cite{HolmgrenRothblum2024} gives $O(L+tq\poly(\log L))$ size for the same
local queries; our bound saves a factor $\log L$ and controls the
leading coefficient.

\paragraph{Relation to Hirahara's reduction.}
In Hirahara's NP-hardness proof for partial MCSP
\cite[Lemmas~8.1--8.3]{Hirahara2022}, locally amplified random tables mask
secret shares.  A satisfying assignment selects tables that let a
completeness circuit recover the secret; Uhlig's theorem keeps this
circuit small despite the amplifier's many queries.  For a source table
of length $L$, our theorem expands the available query range from
$L^{o(1/\log\log L)}$ to $L^\gamma$ for each fixed $\gamma<1$, retaining
size $O_\gamma(L/\log L)$ and depth $O_\gamma(\log L)$ for the table
evaluations.  This strengthens the local synthesis step.  It does not
improve the final approximation exponent in Hirahara's Theorem~8.5:
his choice $\Delta=(\log v)^{1/2}$, with $v$ the number of source variables,
already fits Uhlig's range, and the explicit output length
$2^{O(\log v+\Delta^2)}$ must also remain polynomially bounded.

\section{Implications and open problems}\label{sec:discussion}

We record a standard secure-evaluation consequence and examine the leading
coefficient, deterministic construction, and endpoint copy rates.

\paragraph{Secure batch evaluation.}
Fix integers $m\ge4$ and $0\le c<m/3$, and a synchronous network of $m$
parties with private authenticated channels and broadcast.
For every public $f:\bits^n\to\bits$,
fixed $0\le\gamma<1$, and $1\le t\le2^{\gamma n}$, there is a nonuniform
protocol computing $f^{\times t}$ with perfect security against a static
malicious adversary corrupting at most $c$ parties, using
$O_{\gamma,m}(2^n/n)$ communication bits and $O_{\gamma,m}(n)$ rounds.
The inputs may be distributed arbitrarily, and all parties receive the
output vector.  Apply the Ben-Or--Goldwasser--Wigderson framework
\cite{BenOrGoldwasserWigderson1988}, with the quantitative bound of
\cite[Theorem~1.2 of the full version]{AbrahamAsharovYanai2021}, to
\cref{thm:linear-depth}: secure evaluation of a Boolean circuit with size
$S$ and depth $D$ costs $O_m(S+nt)$ bits and $O_m(D+1)$ rounds, including
input sharing and output reconstruction, and $nt=o_\gamma(2^n/n)$.
The circuit and its menus are public nonuniform protocol data.

\subsection{Lower bounds and the leading coefficient}\label{sec:lower-bounds}

For fixed $0<\gamma<1$ and $t=\lfloor2^{\gamma n}\rfloor$, the normalized
worst-case size has lower limiting coefficient at least $1$ by one-copy
counting, while \cref{thm:coefficient} gives upper limiting coefficient
$1/(1-\gamma)$.  Achieving coefficient $1$ at a positive fixed exponential
copy rate remains unresolved.  It is also open whether the upper
coefficient $1/(1-\gamma)+o_\gamma(1)$ can be attained at depth
$O_\gamma(n)$, as the proof of \cref{thm:linear-depth} allows constant losses.

\Cref{thm:module-bound} supplies a matching lower coefficient for a
restricted resource architecture.  At this batch size, suppose a circuit
is partitioned into at most $2^{\mu n}$ modules, for a fixed $0<\mu<1$,
each with polynomially many input and output ports, and an exterior of
$W$ gates.  If gate-count caps $S=O(2^n/n)$ and $W$ suffice to compute
$f^{\times t}$ for every
$f$, then
\[
 (1-\gamma)S+\gamma W\ge(1-o(1))\frac{2^n}{n}.
\]
Consequently $W=o(2^n/n)$ forces the coefficient $1/(1-\gamma)$.
The modules and their wiring may depend on $f$.  The bound applies to
separately synthesized resource banks of the kind used here, but leaves
unrestricted circuits open.  Its proof in \cref{app:module-bound} cuts the
largest modules at their output ports and counts the remaining circuitry.

\subsection{Open problems}

\paragraph{Can the menus be constructed deterministically?}
\Cref{cor:randomized-construction} constructs circuits attaining the
simultaneous size and depth bounds in time polynomial in the supplied
truth-table length, with high probability of correctness on all batches.
It samples fixed menus without certifying their universal property.
Efficient deterministic construction of such menus, or efficient
certification of a sampled menu, remains open here.

\paragraph{Where does the plateau end?}
The theorem reaches $2^{(1-\varepsilon)n}$ copies for every fixed
$\varepsilon>0$, but does not treat a rate that approaches one with $n$.
Some degradation in that regime is forced.  Suppose $n\ge2$ and $f$ depends
on all $n$ of its variables.  Restrict a circuit for $f^{\times t}$ to the
union of its output cones, and write $s$ for the remaining gate count.
Every one of the $tn$ inputs remains.  Each connected component of the
underlying undirected graph contains an output, so there are at most $t$
components.  The graph therefore has at least $tn+s-t$ edges, while fan-in
at most two permits at most $2s$ edges.  Hence
\[
  \C(f^{\times t})\ge t(n-1).
\]
A bound of the form $A\,2^n/n$ therefore cannot hold once
$t>A\,2^n/(n(n-1))$; in particular, a uniform $O(2^n/n)$ plateau must fail for
$t=\omega(2^n/n^2)$.  The theorem still leaves the detailed mass-production
curve throughout the near-endpoint regime $t=2^{n-o(n)}$ unresolved.
Analyzing this regime requires uniform control of the code rate and circuit
overhead as the parameters vary with $n$.

\paragraph{What happens near the quantum endpoint?}
For $t\le2^n$, \cref{prop:quantum-lower} gives
$M(n,t)=\Omega(nt+(2^n/n)\log_2(t+1))$.
\Cref{thm:quantum} matches this order for each fixed rate below one,
where $nt$ is negligible, but does not settle $t=2^{n-o(n)}$.
In particular, the $\Theta(2^n)$ cost at fixed positive rates cannot
continue to $t=2^n$, where touching all output wires already costs
$\Omega(n2^n)$.  Matching the two-term lower bound near this endpoint,
or optimizing scratch space and $T$ count simultaneously with total
gate count, remains open here.  The state theorem also does not settle
the corresponding optimal tradeoff for arbitrary unitary operations.

\clearpage
\appendix
\section{An explicit recursive construction}\label[appendix]{sec:induction}

The explicitly constructible greedy scheduler gives a recursive proof of
\cref{thm:main}.  The leading-coefficient bound in \cref{thm:coefficient}
uses the menu scheduler from the main text.
Apply the composition bound recursively.  At level $k$, split the input
into $k+1$ approximately equal blocks, encode one block, and invoke the
preceding level on the remaining $k$.

\subsection{Greedy line scheduling}

The geometric count says that a suitable line exists, but a circuit must find
one from the target points.  We do this greedily.  Previously occupied points
rule out directions; projective ranking and sorting let an ordinary Boolean
circuit choose the first direction that remains.

Consider a multiset of $g$ requested information points
$u_1,\ldots,u_g\in A^\ell$.  Process the requests in order.  Suppose disjoint
recovery sets have been selected for requests $1,\ldots,j-1$, and let $U$ be
their union.  Then
\[
  |U|=(j-1)(q-1)<gq.
\]
For the target $u_j$, every point $y\in U\setminus\{u_j\}$ forbids at most one
projective direction, namely the direction of $y-u_j$.  The point $u_j$
itself forbids none because it is omitted from the new recovery set.

\begin{lemma}[Projective indexing]\label{lem:projective-indexing}
Projective directions in $\F_q^\ell$ have rank and unrank circuits of size
$\poly_\ell(b)$, with ranks in $[D_q]_0$.
\end{lemma}

\emph{Implementation idea.}
The indexing normalizes the first nonzero coordinate to one and orders the
normalized vectors by the position of that coordinate and then by their
remaining base-$q$ digits.  The explicit circuits are given in
\cref{app:projective-indexing}.

\begin{lemma}[Greedy line scheduler]\label{lem:scheduler}
If
\begin{equation}\label{eq:group-condition}
  gq<D_q,
\end{equation}
then every multiset of $g$ target points admits pairwise disjoint affine-line
recovery sets.  Given the targets, the recovery sets can be selected by a
Boolean circuit of size
\[
  \widetilde O_\ell(g^2q).
\]
\end{lemma}

\emph{Proof idea.}
At stage $j$, the blocking property in \cref{prop:local-gadget} leaves a valid
direction because fewer than $D_q$ directions are forbidden.  The circuit
ranks the forbidden directions, sorts those ranks, and selects the least
missing one; stage $j$ costs $\widetilde O_\ell(jq)$.  The treatment of
duplicate targets, sentinel ranks, and the least-missing-rank circuit is given
in \cref{app:scheduler-proof}.  Summing over the stages gives the stated cost.

\subsection{Grouping and composition}

For $r\ge1$, define
\[
  L_r(n)=\max_{f:\bits^n\to\bits}\C(f^{\times r}).
\]
For each $f$, choose a circuit for $f^{\times r}$ of size at most $L_r(n)$.  If
a construction supplies fewer than $r$ live inputs, the unused inputs of this
circuit may be fixed arbitrarily and their outputs ignored.  Thus a bound for
$L_r(n)$ also applies to every smaller number of copies.

Partition the $t$ requests into $C$ fixed groups, each of size at most
\[
  g=\left\lceil\frac{t}{C}\right\rceil.
\]
Apply \cref{lem:scheduler} independently inside each group.  A code coordinate
is used at most once per group and therefore at most $C$ times overall.  The
inequalities $\lceil t/C\rceil\le t/C+1$ and
$2t\le t^2/C+C$, the latter being the arithmetic--geometric mean inequality,
show that the total scheduling cost is
\begin{equation}\label{eq:scheduling-cost}
  C\cdot\widetilde O_\ell(g^2q)
  =\widetilde O_\ell\!\left(\frac{t^2q}{C}+Cq\right).
\end{equation}
The construction is valid whenever \eqref{eq:group-condition} holds for
$g=\lceil t/C\rceil$, that is, whenever
\begin{equation}\label{eq:scheduling-condition}
  \left\lceil\frac{t}{C}\right\rceil q<D_q.
\end{equation}

Increasing $C$ reduces scheduling cost and increases the number of suffixes
each resource must handle.  The recursion chooses $C$ to balance these costs.

\paragraph{Bounded-demand invariant.}
Fix $f$ and the local-recovery gadget above.  Each $H_{y,r}$ is one fixed
Boolean function on $d$ variables.  Inside each request group, disjoint
recovery ensures that coordinate $y$ occurs at most once; hence the pair
$(y,r)$ receives at most one suffix from each of the $C$ groups.  Place those
suffixes in the canonical group slots and pad unused slots with a fixed dummy
suffix.  One circuit for $H_{y,r}^{\times C}$, of size at most $L_C(d)$,
serves exactly this bounded demand.  There are $q^\ell b$ such pairs, and no
sharing between different pairs is assumed.

\begin{proposition}[Grouped composition]
\label{prop:grouped-composition}
Fix $\ell\ge2$ and the local-recovery gadget for a split $n=p+d$.  Let
$C\ge1$, set $g=\lceil t/C\rceil$, and suppose that, for every group size
$1\le h\le g$, a Boolean circuit of size at most $S(g,q)$ maps every multiset
of $h$ information points to enumerated pairwise disjoint affine-line recovery
sets.  Then
every $f:\bits^n\to\bits$ satisfies
\begin{equation}\label{eq:grouped-composition}
\C(f^{\times t})
\le \underbrace{q^\ell b\,L_C(d)}_{\text{resource circuits}}
   +\underbrace{C S(g,q)}_{\text{schedule selection}}
   +\underbrace{\widetilde O_\ell\!\left(tq+Cq^\ell\right)}_{
     \text{routing and decoding}}.
\end{equation}
\end{proposition}

\begin{proof}
Apply the assembly from \cref{prop:scheduler-sensitive}, with one scheduler
per group at total cost $CS(g,q)$.
Use \cref{lem:routing} with $M=Cq^\ell$ keys $(\text{group},y)$.
After scatter, fixed wiring transposes the canonical key order to
$(y,\text{group})$, placing all suffixes for one resource together.

The bounded-demand invariant permits evaluating $H_{y,r}^{\times C}$ in
each resource position, at total cost $q^\ell bL_C(d)$.  Gather and decode
as in \cref{prop:scheduler-sensitive}.  Prefix mapping, routing, and decoding
together cost $\widetilde O_\ell(tq+Cq^\ell)$.
\end{proof}

\begin{samepage}
\begin{proposition}[Composition bound]\label{prop:master}
Fix $\ell\ge2$.  Let $n=p+d$ with $p$ sufficiently large depending on
$\ell$, choose $q=2^b$ minimally so that $s_0^\ell b\ge2^p$, and suppose
\cref{eq:scheduling-condition} holds.  Then every
$f:\bits^n\to\bits$ satisfies
\begin{equation}\label{eq:master}
\C(f^{\times t})
\le
\underbrace{q^\ell b\,L_C(d)}_{\text{recursive resource evaluation}}
+\widetilde O_\ell\!\left(
  \underbrace{\frac{t^2q}{C}}_{\text{line scheduling}}
  +\underbrace{tq}_{\text{incidences and decoding}}
  +\underbrace{Cq^\ell}_{\text{resource slots}}
\right).
\end{equation}
\end{proposition}
\end{samepage}

The induction keeps the resource term at scale $2^n/n$ while making all
three overhead terms exponentially smaller.

\begin{proof}
Put $g=\lceil t/C\rceil$.  The scheduling condition lets us apply
\cref{lem:scheduler} with
$S(g,q)=\widetilde O_\ell(g^2q)$.  Moreover,
\[
  Cg^2q
  \le C\left(\frac{t}{C}+1\right)^2q
  =O\!\left(\frac{t^2q}{C}+tq+Cq\right).
\]
Apply \cref{prop:grouped-composition}; the $tq$ term is already present, and
$Cq$ is absorbed by $Cq^\ell$ because $q\ge2$ and $\ell\ge2$.
\end{proof}

\subsection{The induction invariant}

For an integer $k\ge1$, let $\mathsf P_k$ denote the assertion that, for every
fixed $0\le\gamma<k/(k+1)$, there is a constant $A_{k,\gamma}$ such that, for
all sufficiently large $n$, every $f:\bits^n\to\bits$, and every integer
$1\le t\le2^{\gamma n}$,
\[
  \C(f^{\times t})\le A_{k,\gamma}\frac{2^n}{n}.
\]

Every comparison below has a fixed positive linear margin in its base-two
logarithm.  We use without further comment that such a margin absorbs the
$o(n)$ terms, every fixed polynomial factor hidden by $\widetilde O$, and the
$O(1)$ changes caused by floors and ceilings.

\subsection{Two equal blocks: the base case}

\begin{lemma}[Two-block base]\label{lem:base}
$\mathsf P_1$ holds.
\end{lemma}

Here $C=1$, so no mass-production hypothesis is used inside a resource.  The
two equal blocks make the resource term proportional to $2^{n/2}L(n/2)$; the
one-half threshold comes from the quadratic scheduler term $t^2q$.

\begin{proof}
Fix $0\le\gamma<1/2$, let $\eta=1/2-\gamma>0$, and suppose
$t\le2^{\gamma n}$.  Put
\[
  m=\left\lfloor\frac n2\right\rfloor,
  \qquad p=m,
  \qquad d=n-m,
  \qquad C=1.
\]
Thus $n=2m+O(1)$ and all recovery sets will be globally disjoint.  In block
units, the request and code parameters are
\[
  \log_2t\le m-\eta n+O(1),
  \qquad
  \log_2q=\frac m\ell+o(n),
  \qquad
  \log_2D_q=\left(1-\frac1\ell\right)m+o(n).
\]

Choose a sufficiently large fixed $\ell$ so that
\begin{equation}\label{eq:base-margin}
  \frac1\ell<\eta.
\end{equation}
Since $2m/\ell\le n/\ell<\eta n$, this gives
\[
  m-\eta n+\frac m\ell
  <m-\frac m\ell.
\]
The strict linear margin now implies $tq<D_q$ for all sufficiently large $n$.

The resource term in \cref{eq:master} is, by Shannon--Lupanov synthesis,
\[
  q^\ell bL(d)
  =O_\ell\!\left(2^p\frac{2^d}{d}\right)
  =O_\gamma(2^n/n).
\]
The scheduler exponent is
\[
  \log_2(t^2q)
  \le 2m-2\eta n+\frac m\ell+o(n)<n,
\]
where the strict inequality follows from \eqref{eq:base-margin}.  Likewise,
$\log_2(tq)\le m-\eta n+m/\ell+o(n)<n$ and
$\log_2q^\ell=m+o(n)<n$.  After the suppressed polynomial factors are
restored, every overhead term is $2^{(1-\Omega_\gamma(1))n}$ and is negligible
compared with $2^n/n$.  The same strict margins absorb all rounding.
\end{proof}

\subsection{Adding one equal block}

\begin{lemma}[Equal-block step]\label{lem:block-step}
For every integer $k\ge2$, $\mathsf P_{k-1}$ implies $\mathsf P_k$.
\end{lemma}

\begin{proof}
Fix $0\le\gamma<k/(k+1)$.  If $\gamma<(k-1)/k$, the assertion is already
$\mathsf P_{k-1}$.  We may therefore assume
\[
  \frac{k-1}{k}\le\gamma<\frac{k}{k+1}.
\]
Let
\[
  \eta=\frac{k}{k+1}-\gamma>0,
  \qquad
  m=\left\lfloor\frac{n}{k+1}\right\rfloor,
  \qquad
  p=m,
  \qquad
  d=n-m.
\]
Write $n=(k+1)m+r$ with $0\le r\le k$, and choose the number of request
groups to be
\begin{equation}\label{eq:block-parameters}
  C=\left\lceil
       2^{(k-1)m-\eta n/2}
     \right\rceil.
\end{equation}
This choice splits the available slack in half.  Dividing the requests among
$C$ groups puts each group $\eta n/2$ below the one-block exponent $m$;
measured on the remaining $d$ bits, the exponent of $C$ is also a fixed
amount below the threshold supplied by $\mathsf P_{k-1}$.
The assumption $\gamma\ge(k-1)/k$ ensures that the exponent in
\eqref{eq:block-parameters} is positive for all sufficiently large $n$.

Let $\beta=(k-1)/k$.  Since $d=km+r$,
\[
  \beta d-\log_2C
  =\frac{\eta n}{2}+\frac{k-1}{k}r+O(1).
\]
In particular, $C\le2^{(\beta-\eta/4)d}$ for all sufficiently large $n$.
The exponent $\beta-\eta/4$ is fixed and strictly below the threshold for
$\mathsf P_{k-1}$, so the inductive hypothesis gives
\[
  L_C(d)=O_{k,\gamma}(2^d/d).
\]
By \cref{eq:packing}, the resource term is therefore
\[
  q^\ell bL_C(d)
  =O_{k,\gamma,\ell}\!\left(2^m\frac{2^d}{d}\right)
  =O_{k,\gamma,\ell}(2^n/n).
\]

The remaining estimates are easiest to read in block units.  Since
$\log_2t\le\gamma n=km-\eta n+O_k(1)$, the size
$g=\lceil t/C\rceil$ of a request group satisfies
\[
  \log_2g\le m-\frac{\eta n}{2}+O_k(1).
\]
The code parameters satisfy
\[
  \log_2q=\frac m\ell+o(n),
  \qquad
  \log_2D_q=\left(1-\frac1\ell\right)m+o(n).
\]
Choose a sufficiently large fixed $\ell$ so that
\begin{equation}\label{eq:block-margin}
  \frac{2}{(k+1)\ell}<\frac{\eta}{2}.
\end{equation}
Then $2m/\ell<\eta n/2$, and hence
\[
  m-\frac{\eta n}{2}+\frac m\ell
  <m-\frac m\ell.
\]
Together with \eqref{eq:directions}, this is exactly the strict exponent margin
required for \cref{eq:scheduling-condition}.

Finally, the complete overhead ledger in \cref{eq:master} is
\begin{align*}
  \log_2\!\left(\frac{t^2q}{C}\right)
  &\le (k+1)m-\frac{3\eta n}{2}+\frac m\ell+o(n)<n,\\
  \log_2(tq)
  &\le n-m-\eta n+\frac m\ell+o(n)<n,\\
  \log_2(Cq^\ell)
  &\le n-m-\frac{\eta n}{2}+o(n)<n.
\end{align*}
The first strict inequality uses $(k+1)m\le n$ and
\eqref{eq:block-margin}; the other two use $m=\Theta_k(n)$ and $\ell>1$.
Thus every overhead term is exponentially smaller than $2^n/n$.  This proves
$\mathsf P_k$.
\end{proof}

\begin{proof}[Alternative proof of \cref{thm:main}]
By \cref{lem:base,lem:block-step}, $\mathsf P_k$ holds for every $k\ge1$.
Fix $0\le\gamma<1$ and choose $k$ so that $\gamma<k/(k+1)$.  Then
$\mathsf P_k$ gives the asserted bound beyond a threshold $N_\gamma$.  For the
finitely many integers $1\le n<N_\gamma$, use
$\C(f^{\times t})\le t\C(f)$ and enlarge $A_\gamma$ to dominate
$nt\C(f)/2^n$ over all relevant $n$, $f$, and $t$.  This proves the statement
for every $n\ge1$.
\end{proof}

\section{Circuit implementation details}\label[appendix]{app:implementation}

This appendix records the gate-level details behind the scheduler and routing
lemmas.  The main proof uses only their statements and costs.

\subsection{Projective-direction indexing}\label[appendix]{app:projective-indexing}

\begin{proof}[Proof of \cref{lem:projective-indexing}]
Normalize a nonzero vector so that its first nonzero coordinate is one.  The
normalized directions whose first nonzero coordinate is in position
$h\in\{1,\ldots,\ell\}$ have the form
\[
  (0,\ldots,0,1,w_{h+1},\ldots,w_\ell)
\]
and form a block of size $q^{\ell-h}$.  Let
\[
  B_h=\sum_{r=1}^{h-1}q^{\ell-r},
  \qquad
  \operatorname{val}_q(w_{h+1},\ldots,w_\ell)
   =\sum_{r=h+1}^{\ell}\operatorname{int}(w_r)q^{\ell-r},
\]
where $\operatorname{int}(w_r)\in[q]_0$ is the canonical binary
representation.  The rank is
\[
  B_h+\operatorname{val}_q(w_{h+1},\ldots,w_\ell).
\]
The blocks are consecutive and their total size is $D_q$.

To rank, find the first nonzero coordinate, invert it, normalize the vector,
and concatenate the $b$-bit tail coordinates.  To unrank, compare the rank
with the $\ell$ hardwired block intervals, subtract the appropriate $B_h$,
and split the remainder into base-$q$ digits.  Since $q=2^b$, the last step is
bit slicing.  The field operations and comparisons from the circuit
conventions give size $\poly_\ell(b)$ in both directions.
\end{proof}

\subsection{The least-missing-direction circuit}\label[appendix]{app:scheduler-proof}

\begin{proof}[Proof of \cref{lem:scheduler}]
At step $j$, fewer than $D_q$ directions are forbidden, so a valid direction
exists.  This proves the combinatorial statement.

For the circuit bound, represent a projective direction by normalizing its
first nonzero coordinate to one and use the ranking from
\cref{lem:projective-indexing}.  At stage $j$, for each $y\in U$, output the
sentinel rank $D_q$ when $y=u_j$ and otherwise output the projective rank of
$y-u_j$.  At the gate level, first multiplex a fixed nonzero vector in place
of $y-u_j$ when $y=u_j$, rank that always-nonzero vector, and then overwrite
the result by $D_q$ on the equality branch.  Use
$\lceil\log_2(D_q+1)\rceil$ bits per rank.  There are
$(j-1)(q-1)=O(jq)$ such entries.

Sort these ranks as $r_1\le\cdots\le r_s$.  The least missing rank can be
found without a dense $D_q$-bit table.  Rank zero is a candidate if $s=0$ or
$r_1>0$; after position $i$, rank $r_i+1$ is a candidate if
\[
  r_i<D_q-1
  \quad\text{and}\quad
  \bigl(i=s\text{ or }r_{i+1}>r_i+1\bigr).
\]
These tests are applied to the full sorted list, including duplicates and
sentinels.  Only the last occurrence of a non-sentinel rank can generate its
successor, and a sentinel cannot generate a candidate.

\begin{samepage}
A left-to-right
prefix-OR circuit identifies the first candidate in this order.  Multiplexers
select its rank using $O(s\log D_q)$ gates.
\end{samepage}

Since
$|U\setminus\{u_j\}|<D_q$, fewer than $D_q$ distinct non-sentinel ranks are
forbidden, so a candidate exists.  Sorting, comparisons, and priority
selection therefore cost $\widetilde O_\ell(jq)$ gates.  Unrank the selected
direction using \cref{lem:projective-indexing} and enumerate its $q-1$
non-target line points, at cost $q\poly_\ell(b)$.  Unrolling the stages and
summing over $j\le g$ gives $\widetilde O_\ell(g^2q)$.
\end{proof}

\subsection{Four-pass fixed-wire routing}\label[appendix]{app:routing-proof}

\begin{proof}[Proof of \cref{lem:routing}]
For scatter, create one slot record for each of the $M$ table keys and combine
them with the fewer than $tq$ incidence records.  Preserve a copy of every
incidence record by free fan-out.  The uniqueness hypothesis puts at most one
incidence at each key.  After the first pass in the table, a slot therefore
copies the preceding suffix exactly when a type-and-key equality test succeeds;
otherwise it receives a fixed dummy suffix.  The predecessor of the first
position is defined to be a dummy record, so the boundary case uses the same
guard.  The second pass places the filled slots in canonical key order,
ready for the resource circuits.

For gather, combine the preserved incidences with one value record per key.
The third pass places that value immediately before the unique matching
incidence; a guarded local multiplexer copies it into the incidence.  The
fourth pass
puts the first $t(q-1)$ positions in the fixed request-and-line order required
by the decoders.  Thus absent incidences, the first array position, and all
type boundaries are handled by fixed dummy data and the same guarded update,
not by data-dependent wiring.

This is a constant number of sorting networks on
$O(tq+M)$ records.  Each record has width
$\poly_\ell(n,\log q,\log(t+2),\log(M+2))$, so the standard primitives from the
circuit conventions give the stated bound.
\end{proof}

\paragraph{Relation to multiselection.}
The gather half is also an instance of multiselection.  Equivalently, it is a
partial-function inner join in the sense of Holmgren and Rothblum
\cite[Proposition~25 and Sections~6.1--6.2]{HolmgrenRothblum2024}.  Write
$M$ for the number of value slots and $R=t(q-1)$ for the number of
incidences.  Uniqueness of incidence keys gives $R\le M$.  For the resource
tables used here, a key's field coordinates and any group or code number
determine its slot index by fixed arithmetic.  Each incidence in canonical
request--line-position order supplies this index; applying their binary main theorem independently
to the $b$ bit positions therefore gives size
$O(bM+bR\log^3(M+2))$ \cite[Main theorem]{HolmgrenRothblum2024}, within the
bound of \cref{lem:routing}.  The four-pass construction handles scatter and
gather within a single routing bound.

\section{Formalization and exact finite accounting}\label[appendix]{sec:formal-accounting}

This appendix connects the asymptotic proof in \cref{sec:high-rate} with the
Lean theorem.  It records how the formal construction tracks requests, counts
resource circuits, and handles arbitrary input lengths using exact integer
parameters.

\paragraph{Companion and reproducibility.}
\ifanonymous
The public repository URL is omitted during double-blind review.
\else
At revision
\href{https://github.com/SamuelSchlesinger/algebraic-circuits/tree/8dd82c96f44dbeeaca31f4cc96c687c6d87d1489}{\texttt{8dd82c9}},
the Lean~4.33.1/mathlib~4.33.1 module
\texttt{Algebraic.MassProduction} exposes
\texttt{BlockInduction.exponentialMassProduction} for the explicit proof and
\texttt{Nonuniform.realSharpMassProduction} for the nonuniform proof.
Both endpoints use only Lean's standard axioms, with no proof placeholders.
A small executable Python model of the menu verifier and the high-rate
subcode, with exhaustive checks on tiny parameters, accompanies the
manuscript source at revision
\href{https://github.com/SamuelSchlesinger/boolean-mass-production/tree/1c880d92420ea2709451826074a48945caf9fa67/scripts}{\texttt{1c880d9}}.
\fi

\paragraph{Request identity and inactive slots.}
Each request carries a distinct fixed identifier, even when its target and
suffix coincide with another request's.  The scheduler maintains a
permutation of accepted and pending requests, together with disjoint recovery
lines for the accepted prefix.  Each phase preserves all original payloads.
For regular array sizes, a line uses $q$ scalar slots: the zero scalar is
marked inactive, and its recovered Boolean contribution is fixed to zero.
After XOR recovery, only the identifier and the one-bit answer need to be
sorted back into request order.  Thus restoring output order does not require
another permutation of the full point lists.

\paragraph{Charge exactly the resource bank.}
A single hardwired table indexed by the source prefix can return the code
number, information point, and basis-bit selector together.  Batched lookup
handles repeated prefixes and removes runtime division from this version of
the construction.  The resource key is
$(\text{code number},\text{point},\text{basis bit})$; valid scheduled
incidences have distinct keys.  Evaluate precisely the $R=BNb$ Boolean
resource circuits, once each.  Padding the sorting networks adds inactive
routing records, not resource evaluations.  This distinction avoids a
constant-factor loss in the main term.  If $\widehat L(d)$ bounds the chosen
one-copy synthesis circuits, the full cost is bounded by
\[
  R\widehat L(d)+H,\qquad H\le c_\ell 2^p n^7
\]
in the parameter regime below, with $c_\ell$ independent of $n,f,t$.
For each fixed positive integer $Q$, the actual synthesis circuits eventually
admit $\widehat L(d)=\lfloor(Q+1)2^d/(Qd)\rfloor$.
This checked polynomial envelope is coarser than the optimized pass count
in \cref{thm:linear-scheduler}; both give the required negligible overhead.

\paragraph{Whole field blocks, arbitrary input lengths.}
One exact parameter choice fixes positive integers $h,\ell,G,S$, with
$\ell\le2^h$, and sets
\[
  P=\ell h+1<S,\qquad
  m=\lfloor n/S\rfloor,\qquad b=hm,\qquad p=Pm,\qquad d=n-p.
\]
Choose $G+h+1\le h(\ell-1)$ and $\gamma<G/S$.
For sufficiently large $m$, the scheduler batch $2^{Gm}$ covers the requested
copies, including rounding, and has $512\,2^{Gm}q\le D_q$.
Also $2^{Gm}q\le2^p$.  Put the incomplete input block into the suffix; then
\[
  p+d=n,\qquad \frac dn\ge\frac{S-P}{S}>0.
\]
Integer slopes can make $P/S$ arbitrarily close to $\gamma$ from above
while keeping all these inequalities strict.  All dimensions and slopes
are fixed before $n$ varies.

For $A=2^{\ell h}$ the retained dimension on this subsequence is exactly
$K=A^m-(A-1)^m$.  At a fixed positive integer precision $Q$, the elementary
rate estimate is
\[
  m\ge Q(A-1)\quad\Longrightarrow\quad QA^m\le(Q+1)K.
\]
The formal packing uses $B=\lfloor2^p/(Kb)\rfloor+1$, charging at most one
extra codeword.  Since $Nb=2^{\ell hm}hm$ and $p=(\ell h+1)m$, eventually
$QNb\le2^p$.  Consequently $QR\le(Q+2)2^p$.
The number $B$ need not be polynomial in this alternative parameter choice;
the bounds on the total resource bank and the key widths are what matter.
Finally, $m\le d$ and every fixed polynomial in $n$ is eventually smaller
than $2^m$.  These facts absorb $H$ after normalization by $n/2^n$.

\paragraph{Uniform quantifiers.}
For a rational rate $u/v<1$ and every positive integer precision $J$, the
checked integer endpoint supplies one cutoff, uniform over $f$ and all
positive $t\le2^{\lfloor un/v\rfloor}$, with
\[
  J(v-u)n\,\C(f^{\times t})\le(J+1)v2^n.
\]
The final real-rate theorem chooses a slightly larger rational rate and a
large enough precision inside the requested additive error.  Upward exponent
rounding is covered by the scheduler's strict capacity margin.  The result
has exactly the real-rate and $\varepsilon$ quantifiers of
\cref{thm:coefficient}, without substituting growing dimensions into a
fixed-dimension estimate.

\section{A lower bound for modules with polynomial interfaces}
\label[appendix]{app:module-bound}

We count the functions representable by circuits built from modules
whose input and output interfaces are small.  The internal gates of
different modules are disjoint.  Each module receives all external
signals through its input ports, and its internal gates may be used
outside the module only through its designated output ports.  All other
gates form the \emph{exterior}.  The assembled circuit is acyclic;
wiring, constants, and output designations remain free.

\begin{theorem}[Module counting bound]\label{thm:module-bound}
Fix $0<\gamma,\mu<1$ and polynomial bounds $a(n),b(n)\ge1$.  Put
$N=2^n$ and $t=\lfloor2^{\gamma n}\rfloor$.  Let $\mathcal F_n(S,W)$
be the family of functions $f:\bits^n\to\bits$ for which $f^{\times t}$
has a circuit with at most $S$ gates, partitioned into at most
$2^{\mu n}$ modules and an exterior of at most $W$ gates, with at most
$a(n)$ input and $b(n)$ output ports per module.
For $0\le W\le S=O(N/n)$,
\[
 |\mathcal F_n(S,W)|
 \le 2^{\,n((1-\gamma)S+\gamma W)+o(N)}.
\]
The error is uniform under the fixed rate, port, and size bounds.
In particular, gate-count caps covering every $f$ must satisfy
\begin{equation}\label{eq:module-lower}
 (1-\gamma)S+\gamma W\ge(1-o(1))\frac{N}{n}.
\end{equation}
Any fixed positive gap below $N/n$ on the left makes the fraction
of representable functions exponentially small in $N$.
\end{theorem}

The proof cuts the largest modules at fewer than $t$ output ports.
For each fixed remainder, their possible replacements account for only
$2^{o(N)}$ functions.  The modules left behind are small enough to
describe at a lower cost per gate.  Two elementary counting facts
make these steps precise.

\begin{lemma}[Cutting module outputs]\label{lem:module-cut}
Fix the exterior, all interconnections, and all but a selected group of
modules whose output ports total $r<t$.  Let $\mathcal F$ be the family
of functions $f$ for which some choice of the selected modules completes
the circuit to $f^{\times t}$.  Then
\[
 |\mathcal F|\le\sum_{i=0}^r\binom Ni\le(r+1)N^r.
\]
\end{lemma}

\begin{proof}
Regard a Boolean function as a subset of its $N$-point domain.
A set of points is \emph{shattered} if the family realizes every Boolean
labeling of those points.  If $\mathcal F$ shattered $r+1$ distinct
points, place them in the first
$r+1$ request slots of a fixed batch and fill the remaining slots
arbitrarily.  Cut the selected output ports and treat their values as
$r$ independent Boolean sources.  There are at most $2^r$ resulting
output vectors from the fixed remainder, whereas shattering requires
$2^{r+1}$ labelings of those points.  This is impossible.  Cutting all
selected outputs also handles connections between selected modules.

Thus the VC dimension of $\mathcal F$, the largest size of a shattered
set, is at most $r$.  Sauer's bound
\cite{Sauer1972} gives the first inequality.  For completeness, its
induction splits a family by the last domain point, with projected
families $\mathcal F_0,\mathcal F_1$.  The original cardinality is
$|\mathcal F_0\cup\mathcal F_1|+|\mathcal F_0\cap\mathcal F_1|$.
The union has VC dimension at most $r$ and the intersection at most
$r-1$, so Pascal's identity gives the binomial sum.  The empty-family
and zero-dimension cases supply the induction bases.  Finally
$\binom Ni\le N^i$ proves the second inequality.
\end{proof}

\begin{lemma}[Descriptions with several outputs]\label{lem:multi-output-count}
After deleting unused gates, circuits with $A$ named input sources,
$Q$ designated outputs, and $u$ gates can be described, up to relabeling
of internal gates, in at most $12^u2^Q(u+A+2)^{u+Q}$ ways.
The base-two logarithm of this count is at most
\[
 u\log_2 12+Q+(u+Q)\log_2(u+A+2).
\]
\end{lemma}

\begin{proof}
Traverse the output cones depth first, with the roots in output order.
At each root or incoming edge, record whether a new gate is discovered
or an existing gate, input, or constant is named.  There are at most
$2u+Q$ such flags.  Exactly $u$ positions discover a gate, whose type
has three possibilities, leaving at most $u+Q$ names from a set of
size $u+A+2$.  Number gates in discovery order.  Padding unused
description positions gives the stated count; invalid descriptions
can only increase it.
\end{proof}

\begin{proof}[Proof of \cref{thm:module-bound}]
First fix exact gate counts $S,W$ and a module count $M\le2^{\mu n}$.
Pad the ports to the bounds $a=a(n)$ and $b=b(n)$, rounding these
bounds up to integers if needed.  Put $B=S-W$ and
\[
 k=\left\lfloor\frac{t-1}{b}\right\rfloor.
\]
Remove the $k$ largest modules by gate count, or all modules if there
are fewer.  They expose at most $kb<t$ output ports.  For each fixed
remainder, \cref{lem:module-cut} bounds the logarithm of the number of
possible functions by
\[
 \log_2(r+1)+r\log_2N=O(tn)=o(N).
\]
This step accounts for the removed modules without describing their
potentially large interiors.

Every remaining module has at most $B/(k+1)$ gates.  Applying
\cref{lem:multi-output-count} to their $a$ named inputs and $b$ outputs
shows that the logarithm of their total description count is at most
\[
 B\log_2\!\left(\frac{B}{k+1}+a+2\right)
 +O(B)+O\bigl(Mb\log(S+a+2)\bigr).
\]
Here $\log_2(k+1)=\gamma n+O(\log(n+2))$.  Since $a,b$ are polynomial
and $\mu<1$, the displayed bound is at most
\[
 (1-\gamma)nB+o(N).
\]
This remains true if no modules remain.  Specifying module gate counts
and the removed subset costs $O(M\log(S+1)+M)=o(N)$ additional bits.

To include all interconnections, describe the exterior as a circuit
on the $nt$ original inputs and the $Mb$ module output ports, with
designated outputs supplying the $Ma$ module input ports and the $t$
final outputs.  In \cref{lem:multi-output-count}, take
\[
 A_{\rm ext}=nt+Mb,\qquad Q_{\rm ext}=Ma+t,\qquad u=W.
\]
This description includes direct wires between modules, constant
connections, and output designations.  Because $\mu,\gamma<1$,
the terms involving $Q_{\rm ext}$ are $o(N)$, while
$\log_2(W+A_{\rm ext}+2)\le n+O(\log(n+2))$.  The exterior therefore
contributes at most $nW+o(N)$ to the logarithm.  Descriptions that create cycles
when the modules are reattached merely overcount the valid circuits.

For each such description, the cut bound supplies at most $2^{o(N)}$
functions.  Multiplication of the counts gives
\[
 \log_2|\mathcal F_n(S,W)|
 \le (1-\gamma)n(S-W)+nW+o(N)
 =n((1-\gamma)S+\gamma W)+o(N).
\]
Summing over exact counts up to the caps adds only $o(N)$ to the
logarithm; the final expression is increasing in both caps.
There are $2^N$ truth tables, proving \eqref{eq:module-lower} and the
fixed-gap assertion.
\end{proof}

When $W=o(N/n)$, the required coefficient of $S$ is at least
$1/(1-\gamma)$.  Conversely, if $S\le(1+o(1))N/n$ suffices for every
function, \eqref{eq:module-lower} forces $W\ge(1-o(1))N/n$ and hence
$S-W=o(N/n)$.  Thus coefficient one cannot be attained by keeping a
positive fraction of the leading cost inside such a bank of modules.
Substantial sharing across module interiors, larger interfaces, or a
leading-order exterior can escape these hypotheses; the theorem makes
no unrestricted lower-bound claim.

\ifanonymous\else
\section*{Acknowledgements}

I thank Shreyas Srinivas for helpful feedback on the exposition and the
quantum literature.  I am especially grateful for his suggestion to attack
the quantum variant of the mass-production problem.
\fi

\section*{AI Disclosure}

This project began in an interactive conversation on ChatGPT.com that was
exported to the author's computer as a \LaTeX{} document.  Later interactive
sessions used Claude Fable~5 and 5.1 through Claude Code and ChatGPT~5.6 Sol
through Codex.  These systems materially influenced the mathematical
experiments, broad literature searches, proof strategy, and internal
constructions, including the scheduling and recursive-composition approach;
their role went beyond prose editing.  The author selected the directions to
pursue, checked and revised the proofs and literature claims, and takes full
responsibility for the correctness, originality, and integrity of the paper.
Codex also assisted with development and testing of the nonuniform menu
scheduler and the high-rate-code refinement, and with their Lean formalization
through the sharp real-rate theorem.  Codex also assisted with the
quantum state preparation application and its packing lower bound,
and with the local hardness amplification application.

\clearpage
\begin{thebibliography}{HPPV20}
\small
\interlinepenalty=10000

\bibitem[AAY21]{AbrahamAsharovYanai2021}
Ittai Abraham, Gilad Asharov, and Avishay Yanai.
\newblock Efficient perfectly secure computation with optimal resilience.
\newblock In \emph{Theory of Cryptography---TCC 2021, Part II}, LNCS 13043,
pages 66--96. Springer, 2021.
\newblock \href{https://doi.org/10.1007/978-3-030-90453-1_3}{doi:10.1007/978-3-030-90453-1\_3}.
\newblock Full version: \url{https://eprint.iacr.org/2021/1206}.

\bibitem[Alb89]{Albrecht1989}
Andreas Albrecht.
\newblock On simultaneous realizations of Boolean functions, with applications.
\newblock In \emph{Parcella '88: Proceedings of the Fourth International
Workshop on Parallel Processing by Cellular Automata and Arrays}, LNCS 342,
pages 51--56. Springer, 1989.
\newblock \href{https://doi.org/10.1007/3-540-50647-0_102}{doi:10.1007/3-540-50647-0\_102}.

\bibitem[Bat68]{Batcher1968}
Kenneth E. Batcher.
\newblock Sorting networks and their applications.
\newblock In \emph{Proceedings of the 1968 Spring Joint Computer Conference},
AFIPS Conference Proceedings 32, pages 307--314, 1968.
\newblock \href{https://doi.org/10.1145/1468075.1468121}{doi:10.1145/1468075.1468121}.

\bibitem[BGW88]{BenOrGoldwasserWigderson1988}
Michael Ben-Or, Shafi Goldwasser, and Avi Wigderson.
\newblock Completeness theorems for non-cryptographic fault-tolerant
distributed computation.
\newblock In \emph{Proceedings of the 20th Annual ACM Symposium on Theory
of Computing}, pages 1--10. ACM, 1988.
\newblock \href{https://doi.org/10.1145/62212.62213}{doi:10.1145/62212.62213}.

\bibitem[DN06]{DawsonNielsen2006}
Christopher M. Dawson and Michael A. Nielsen.
\newblock The Solovay--Kitaev algorithm.
\newblock \emph{Quantum Information \& Computation}, 6(1):81--95, 2006.
\newblock \url{https://arxiv.org/abs/quant-ph/0505030}.

\bibitem[FM05]{FrandsenMiltersen2005}
Gudmund Skovbjerg Frandsen and Peter Bro Miltersen.
\newblock Reviewing bounds on the circuit size of the hardest functions.
\newblock \emph{Information Processing Letters}, 95(2):354--357, 2005.
\newblock \href{https://doi.org/10.1016/j.ipl.2005.03.009}{doi:10.1016/j.ipl.2005.03.009}.

\bibitem[FVY15]{FazeliVardyYaakobi2015}
Arman Fazeli, Alexander Vardy, and Eitan Yaakobi.
\newblock PIR with low storage overhead: Coding instead of replication.
\newblock arXiv:1505.06241, 2015.
\newblock \url{https://arxiv.org/abs/1505.06241}.

\bibitem[GKS13]{GuoKoppartySudan2013}
Alan Guo, Swastik Kopparty, and Madhu Sudan.
\newblock New affine-invariant codes from lifting.
\newblock In \emph{Proceedings of the 4th Innovations in Theoretical Computer
Science Conference}, pages 529--539. ACM, 2013.
\newblock \href{https://doi.org/10.1145/2422436.2422494}{doi:10.1145/2422436.2422494}.
\newblock \url{https://arxiv.org/abs/1208.5413}.

\bibitem[GKW26]{GossetKothariWu2026}
David Gosset, Robin Kothari, and Kewen Wu.
\newblock Quantum state preparation with optimal $T$-count.
\newblock \emph{Quantum}, 10:2168, 2026.
\newblock \href{https://doi.org/10.22331/q-2026-07-22-2168}{doi:10.22331/q-2026-07-22-2168}.
\newblock \url{https://arxiv.org/abs/2411.04790}.

\bibitem[HKW25]{HugginsKhattarWiebe2025}
William J. Huggins, Tanuj Khattar, and Nathan Wiebe.
\newblock Productionizing quantum mass production.
\newblock arXiv:2506.00132, 2025.
\newblock \url{https://arxiv.org/abs/2506.00132}.

\bibitem[HP95]{HiltgenPaterson1995}
Alain P. Hiltgen and Mike S. Paterson.
\newblock $\mathrm{PI}_k$ mass production and an optimal circuit for the
Ne\v{c}iporuk slice.
\newblock \emph{Computational Complexity}, 5(2):132--154, 1995.
\newblock \href{https://doi.org/10.1007/BF01268142}{doi:10.1007/BF01268142}.

\bibitem[Hir22]{Hirahara2022}
Shuichi Hirahara.
\newblock NP-hardness of learning programs and partial MCSP.
\newblock In \emph{2022 IEEE 63rd Annual Symposium on Foundations of
Computer Science (FOCS)}, pages 968--979. IEEE, 2022.
\newblock \href{https://doi.org/10.1109/FOCS54457.2022.00095}{doi:10.1109/FOCS54457.2022.00095}.
\newblock Full version, used for theorem numbering:
\url{https://eccc.weizmann.ac.il/report/2022/119/}.

\bibitem[HPPV20]{HolzbaurEtAl2020}
Lukas Holzbaur, Rina Polyanskaya, Nikita Polyanskii, and Ilya Vorobyev.
\newblock Lifted Reed--Solomon codes with application to batch codes.
\newblock In \emph{2020 IEEE International Symposium on Information Theory},
pages 634--639. IEEE, 2020.
\newblock \href{https://doi.org/10.1109/ISIT44484.2020.9174402}{doi:10.1109/ISIT44484.2020.9174402}.
\newblock arXiv:2001.11981.

\bibitem[HR24]{HolmgrenRothblum2024}
Justin Holmgren and Ron D. Rothblum.
\newblock Linear-size Boolean circuits for multiselection.
\newblock In \emph{39th Computational Complexity Conference (CCC 2024)},
LIPIcs 300, Article 11, pages 11:1--11:20, 2024.
\newblock \href{https://doi.org/10.4230/LIPIcs.CCC.2024.11}{doi:10.4230/LIPIcs.CCC.2024.11}.

\bibitem[IW97]{ImpagliazzoWigderson1997}
Russell Impagliazzo and Avi Wigderson.
\newblock $\mathrm{P}=\mathrm{BPP}$ if $\mathrm{E}$ requires exponential
circuits: Derandomizing the XOR lemma.
\newblock In \emph{Proceedings of the 29th Annual ACM Symposium on Theory
of Computing}, pages 220--229. ACM, 1997.
\newblock \href{https://doi.org/10.1145/258533.258590}{doi:10.1145/258533.258590}.

\bibitem[IKOS04]{IKOS2004}
Yuval Ishai, Eyal Kushilevitz, Rafail Ostrovsky, and Amit Sahai.
\newblock Batch codes and their applications.
\newblock In \emph{Proceedings of the 36th Annual ACM Symposium on Theory of
Computing}, pages 262--271. ACM, 2004.
\newblock \href{https://doi.org/10.1145/1007352.1007396}{doi:10.1145/1007352.1007396}.

\bibitem[Kre23]{Kretschmer2023}
William Kretschmer.
\newblock Quantum mass production theorems.
\newblock In \emph{18th Conference on the Theory of Quantum Computation,
Communication and Cryptography}, LIPIcs 266, Article 10, pages 10:1--10:11,
2023.
\newblock \href{https://doi.org/10.4230/LIPIcs.TQC.2023.10}{doi:10.4230/LIPIcs.TQC.2023.10}.

\bibitem[Lup58]{Lupanov1958}
Oleg B. Lupanov.
\newblock On a method of circuit synthesis.
\newblock \emph{Izvestiya VUZ, Radiofizika}, 1(1):120--140, 1958.
\newblock In Russian.
\newblock \url{https://radiophysics.unn.ru/issues/1958/1/120}.

\bibitem[PV19]{PolyanskiiVorobyev2019}
Nikita Polyanskii and Ilya Vorobyev.
\newblock Constructions of batch codes via finite geometry.
\newblock In \emph{2019 IEEE International Symposium on Information Theory},
pages 360--364. IEEE, 2019.
\newblock \href{https://doi.org/10.1109/ISIT.2019.8849736}{doi:10.1109/ISIT.2019.8849736}.
\newblock arXiv:1901.06741.

\bibitem[Pau76]{Paul1976}
Wolfgang J. Paul.
\newblock Realizing Boolean functions on disjoint sets of variables.
\newblock \emph{Theoretical Computer Science}, 2(3):383--396, 1976.
\newblock \href{https://doi.org/10.1016/0304-3975(76)90089-X}{doi:10.1016/0304-3975(76)90089-X}.

\bibitem[Sau72]{Sauer1972}
Norbert Sauer.
\newblock On the density of families of sets.
\newblock \emph{Journal of Combinatorial Theory, Series A}, 13(1):145--147, 1972.
\newblock \href{https://doi.org/10.1016/0097-3165(72)90019-2}{doi:10.1016/0097-3165(72)90019-2}.

\bibitem[SBM06]{ShendeBullockMarkov2006}
Vivek V. Shende, Stephen S. Bullock, and Igor L. Markov.
\newblock Synthesis of quantum logic circuits.
\newblock \emph{IEEE Transactions on Computer-Aided Design of Integrated
Circuits and Systems}, 25(6):1000--1010, 2006.
\newblock \href{https://doi.org/10.1109/TCAD.2005.855930}{doi:10.1109/TCAD.2005.855930}.
\newblock \url{https://arxiv.org/abs/quant-ph/0406176}.

\bibitem[Sha49]{Shannon1949}
Claude E. Shannon.
\newblock The synthesis of two-terminal switching circuits.
\newblock \emph{Bell System Technical Journal}, 28(1):59--98, 1949.
\newblock \href{https://doi.org/10.1002/j.1538-7305.1949.tb03624.x}{doi:10.1002/j.1538-7305.1949.tb03624.x}.

\bibitem[Sur24]{Suruga2024}
Daiki Suruga.
\newblock Direct sum theorems beyond query complexity.
\newblock arXiv:2408.15570, 2024. Version~2, revised January~2025.
\newblock \url{https://arxiv.org/abs/2408.15570v2}.

\bibitem[Uhl74]{Uhlig1974}
Dietmar Uhlig.
\newblock On the synthesis of self-correcting schemes from functional elements
with a small number of reliable elements.
\newblock \emph{Mathematical Notes of the Academy of Sciences of the USSR},
15(6):558--562, 1974.
\newblock English translation.
\newblock \href{https://doi.org/10.1007/BF01152835}{doi:10.1007/BF01152835}.

\bibitem[Uhl92]{Uhlig1992}
Dietmar Uhlig.
\newblock Networks computing Boolean functions for multiple input values.
\newblock In M. S. Paterson, editor, \emph{Boolean Function Complexity},
London Mathematical Society Lecture Note Series 169, pages 165--173. Cambridge
University Press, 1992.
\newblock \href{https://doi.org/10.1017/CBO9780511526633.013}{doi:10.1017/CBO9780511526633.013}.

\bibitem[Weg87]{Wegener1987}
Ingo Wegener.
\newblock \emph{The Complexity of Boolean Functions}.
\newblock Wiley--Teubner, 1987.
\newblock \url{https://eccc.weizmann.ac.il/static/books/The_Complexity_of_Boolean_Functions/}.

\end{thebibliography}

\end{document}
